% ----------------------
\subsection{Section 76 (16 February 1832)}
% ----------------------
	\label{DC76}
	
	% -----
	\subsubsection{The ``poetic paraphrase'' of DC 76}
	% -----
	
		For the poetic version of this section, likely written by W. W. Phelps, see \hyperref[EXSS-DC76-POETIC]{here}.
	
	% -----
	\subsubsection{Historical Background}
	% -----
		\label{DC76-HB}
		
		% --
		\paragraph{JSP D2}
		% --
		
			``On 16 February 1832, JS and Sidney Rigdon saw a vision `concerning the economy of God and his vast creation throughout all eternity,' likely while in the upstairs bedroom of the John and Alice (Elsa) Johnson home in Hiram, Ohio. This vision came after JS returned from the January conference in Amherst, Ohio, and after he resumed his work of revising the New Testament at the Johnson home, with Rigdon working as scribe. According to a later JS history, revelations JS had dictated up to February 1832 showed `that many important points, touching the salvation of man, had been taken from the Bible, or lost before it was compiled.' Included in these `important points' was information on what happens after death. This information led JS to think `that if God rewarded every one according to the deeds done in the body, the term `heaven,'...must include more kingdoms than one.' According to the description of the vision, on 16 February 1832, JS and Rigdon reviewed John 5:29, wherein Jesus Christ prophesies that the dead will `come forth; they that have done good, unto the resurrection of life; and they that have done evil, unto the resurrection of damnation.' They reported that when they pondered this verse they together beheld a vision of what awaited humankind after death.
			
			``JS and Rigdon's description of the vision outlined three levels of heavenly glory---celestial, terrestrial, and telestial---and the requirements for entrance into each. According to their report, every person who lived on earth---apart from followers of Satan known as `sons of perdition'---would spend the afterlife in one of these `kingdoms.' These concepts differed considerably from views of the afterlife held by most Protestant churches that the souls of the `righteous' are received into heaven while the `wicked' are cast into hell. Other thinkers and theologians, however, had conceptions of heaven that were more similar to JS and Rigdon's vision: The Universalist church, with which JS's grandfather Asael Smith had affiliated, proclaimed that Christ would temporarily punish sinners but eventually redeem all people. Emanuel Swedenborg, a Swedish scientist and mystic, posited in the mid-1700s that heaven consisted of three different levels (celestial, spiritual, and natural). Alexander Campbell, Rigdon's former associate in the Disciples of Christ, also wrote about `three kingdoms'---the Kingdom of Law, the Kingdom of Favor, and the Kingdom of Glory. Campbell's Kingdoms of Law and Favor, however, could be experienced during mortal life, and only the Kingdom of Glory was reserved for the afterlife. In describing these three kingdoms, Campbell wrote that the first was entered through birth, the second through baptism, and the third through good works. One differed from the next, Campbell declared, `as the sun excelled a star.'
			
			``Neither JS nor Rigdon described in detail how the vision occurred---only that they both saw it at the same time. The shared nature of the vision made it somewhat unusual. Although some experiences of angelic visitations had multiple participants---including the appearance of John the Baptist to JS and Oliver Cowdery in May 1829 and the appearance of an angel showing the gold plates to Cowdery, David Whitmer, and Martin Harris in the summer of 1829---most visions and revelations had been experienced by JS alone. Philo Dibble, who claimed to have been with JS and Rigdon in the Johnson home when the vision occurred, later recounted that JS and Rigdon sat in the upstairs room, where they had conducted much of their work on the Bible revision, with twelve other men. By turns, either JS or Rigdon would ask, `What do I see?' and then relate the scene, after which the other would reply, `I see the same.' There is no indication in Dibble's account that anyone was recording the vision as it occurred; instead, Dibble said there was `not a sound nor motion made by anyone' in the room. Dibble recalled that neither JS nor Rigdon `moved a joint or limb during the time I was there.'
			
			``The text of the vision account itself contains JS and Rigdon's descriptions of what they saw, interspersed with the voice of Deity explicating parts of the vision. According to this record, Jesus Christ conversed with JS and Rigdon during the vision. In some ways, such interaction parallels other visions recounted in the Bible and the Book of Mormon. The Book of Mormon, for example, describes a prophet named Nephi being taken by `the spirit' to `an exceeding high mountain,' where he is shown, among other things, Jesus Christ's life on earth. Nephi's account of this vision records both what he saw and questions and statements made by `the spirit' and by an angel as they observed the vision with Nephi.
			
			``How or precisely when JS and Rigdon recorded the experience of seeing the vision is unknown. According to the written account of the vision, JS and Rigdon were commanded four separate times to record what they were seeing. It may be that they made a record after each command and then proceeded. Alternatively, JS and Rigdon may not have recorded anything until after the vision concluded. They were instructed to write their account while they were `yet in the spirit,' and the text in the account indicates that they did so. If JS and Rigdon recorded the event after its conclusion, therefore, they apparently did so soon thereafter. The first part of the vision (beginning `Here O ye heavens' and ending `...yet entered into the heart of man') seems similar in language and style to JS's revelations, suggesting that JS may have dictated the first part separately from the rest of the account. It is unknown whether JS dictated the entire record to Rigdon, Rigdon wrote it himself, or the two worked collaboratively. Whatever the case, both apparently signed the record after its preparation as a testament to its legitimacy. The placement of Rigdon's signature before JS's indicates that Rigdon was serving as the scribe for the account and signed after completing it. Moreover, JS and Rigdon's vision occurred, as their own account attests, `as we sat doing the work of translation'---in which process Rigdon served as a scribe. Rigdon had also recently transcribed some of JS's dictated revelations.
			
			``The earliest extant copy of the account of the vision is in Frederick G. Williams's handwriting, with a few lines penned by JS. Williams was doing scribal work for JS in July 1832 (and possibly as early as February). He and JS copied the account of the vision into a new blank book---probably during February or March 1832---making the vision the first item in Revelation Book 2. Indeed, JS may have purchased the book for the purpose of entering the account of the vision, given the vision's emphasis on writing down what was seen. Additional copies of the document quickly circulated.
			
			``Reaction to the contents of the vision varied among individuals in the Hiram community and among church members in other areas. Some, such as Lincoln Haskins, thought it contained `great \& marvilous things.' Others struggled to reconcile its concepts of the afterlife with traditional notions of heaven and hell. It also became a target of criticism by outsiders, some of whom regarded it as both `pompous' and an attempt to `embrace and teach Universalism.' A later JS history stated that the vision transcended the knowledge of the afterlife available at the time, declaring that `nothing could be more pleasing to the Saint, upon the order of the kingdom of the Lord, than the light which burst upon the world, through the . . . vision.'''
			
		% --
		\paragraph{Philo Dibble, 1877, Payson Ward General Minutes}
		% --
			
			``\uline{Elder} \uline{Philo} \uline{Dibble} of Springville next addressed the Congregation: Said he was a stranger to most of this congregation, but was not a stranger to the gospel, He had been in the church over 47 years. He was among the first 50 that embraced the gospel in this dispensation. He was present when Jos. Smith and Sidney Rigdon when they had that glorious vision of the creation \&c. It was shown unto them that a hundred and forty four thousand should stand on the earth in the last days as Saviors of men. We all agreed in the heavens to come onto this earth and pass through this probation having forgotten all that we knew before. We are here for trial and they who hold out faithful will receive the crown. He was familiar with the Prophet Joseph from boyhood and knew him to be a man of God. Never loved any one as he loved Joseph. The Lord told the prophet when this church was organized with six members, that he was laying the foundation of a great work. It has already grown to greater dimensions than man could have dreamed of. He was present when the heavens were opened and men saw within the veil and saw the man of sin revealed and many other great and glorious things'' (\hyperref[EMS-1877-PD-7-JAN-1877]{Payson Ward general minutes}).
			
		% --
		\paragraph{Philo Dibble, 1882, \textit{Early Scenes in Church History}}
		% --
		
			``When I arrived at Father Johnson's the next morning, Joseph and Sidney had just finished washing up from being tared and feathered the night before. Joseph said to Sidney: `We can now go on our mission to Jackson County' (alluding to a commandment given them while they were translating, but which they concluded not to attend to until they had finished that work). I felt to regret very much that I had not been with them the evening before, but it was perhaps providential that I was not. On a subsequent visit to Hiram, I arrived at Father Johnson's just as Joseph and Sidney were coming out of the vision alluded to in the Book of Doctrine and Covenants, in which mention is made of the three glories. Joseph wore black clothes, but as this time seemed to be dressed in an element of glorious white, and his face shone as if it were transparent, but I did not see the same glory attending Sidney. Joseph appeared as strong as a lion, but Sidney seemed as weak as water, and Joseph, noticing his condition smiled and said, `Brother Sidney is not as used to it as I am''' (\textit{Early Scenes in Church History}, Juvenile Instructor Office, Salt Lake City, 1882, 80--81).
			
		% --
		\paragraph{Philo Dibble, 1892, JI}
		% --
		
			```The vision which is recorded in the Book of Doctrine and Covenants was given at the house of `Father Johnson,' in Hyrum, Ohio, and during the time that Joseph and Sidney were in the spirit and saw the heavens open, there were other men in the room, perhaps twelve, among whom I was one during a part of the time---probably two-thirds of the time,---I saw the glory and felt the power, but did not see the vision.
			
			``The events and conversation, while they were seeing what is written (and many things were seen and related that are not written,) I will relate as minutely as is necessary.
			
			```Joseph would, at intervals, say: `What do I see?' as one might say while looking out the window and beholding what all in the room could not see. Then he would relate what he had seen or what he was looking at. Then Sidney replied, `I see the same.' Presently Sidney would say `what do I see?' and would repeat what he had seen or was seeing, and Joseph would reply, `I see the same.'
			
			```This manner of conversation was repeated at short intervals to the end of the vision, and during the whole time not a word was spoken by any other person. Not a sound nor motion made by anyone but Joseph and Sidney, and it seemed to me that they never moved a joint or limb during the time I was there, which I think was over an hour, and to the end of the vision.
			
			```Joseph sat firmly and calmly all the time in the midst of a magnificent glory, but Sidney sat limp and pale, apparently as limber as a rag, observing which, Joseph remarked, smilingly, `Sidney is not used to it as I am''' (\textit{Juvenile Instructor} 27:303--304 [15 May 1892]).
			
		% --
		\paragraph{Philo Dibble, undated}
		% --
		
			``At another time I had occasion to visit Hyrum. When I arrived at FatherJohnson's Joseph and Sidney were just coming out of the vision of the Three Glories as recorded in the Book of Doctrine and Covenants. Joseph was dressed in black, but he seemed to be enveloped in an element of glorious white. His face shone as if it were transparent. I did not see the same glory attending Sidney. Joseph appeared as strong as a lion, but Sidney seemed as weak as water. Joseph, seeing this situation of Sidney, smiled and said, `Sidney is not as used to it as I am''' (CHL, MS 15447, Philo Dibble reminiscences, undated, PDF 13).
			
		% --
		\paragraph{JS saw more than what is recorded in the vision}
		% -- 
		
			``There are \sout{are} Some things in my own bosom that must remain there. If Paul could say I knew a man who ascended to the third heaven \& saw things unlawful for man to utter, \uline{I} \uline{more}'' (\hyperref[EMS-1843-JS-21-MAY-1843-HC]{21 May 1843}, JSP D12:326).
			
		% --
		\paragraph{Reception of the vision in section 76}
		% --
		
			For discussion about how people received this vision, see HDDC 928--933;
		
	% -----
	\subsubsection{Versions}
	% -----
		\label{DC76-V}
		
		See the \href{https://drive.google.com/file/d/1goAvNz8jS4R8slYfMG_db55Ty2z_vmgK/view?usp=sharing}{sections comparison} document.
	
		\begin{compactdesc}
			\item[JSP] \hyperref[EMS-1832-JS-16-FEB-1832]{16 February 1832}
			\item[EMS] 1:2 (July 1832), 10--11
			\item[BC] ---
			\item[DC 1835] Section 91, 225--231
			\item[MS] 2:2 (June 1841), 17--21
			\item[DC 1844] Section 92, 349--359
			\item[TS] 5:14 (592--595)
			\item[MS] 14:9 (24 April 1852), 129--132
		\end{compactdesc}

	% -----
	\subsubsection{Section 76:1} % *DC76-1 
	% -----
		\label{DC76-1}

			\footnote{%
				\label{DC76-1-economy}%
				In the version in JSP (see it \hyperref[EMS-1832-JS-16-FEB-1832]{here}), the text contains the following preface:
					\begin{quote}
						A vision of Joseph \& Sidney [Rigdon] February 16\supers{th.} 1832 given in Portage County Hiram Township State of Ohio in North Ame[r]ica which they saw concerning the church of the first born and concerning the economy of God and his vast creation througout all eternity
					\end{quote}
				Here are some 1828 definitions for ``economy'':
					\begin{compactitem}
						\item[1] Primarily, the management, regulation and government of a family or the concerns of a household.
						\item[4] The disposition or arrangement of any work; as the \textit{economy} of a poem.
						\item[5] A system of rules, regulations, rites and ceremonies; as the Jewish \textit{economy}.
						\item[6] The regular operation of nature in the generation, nutrition and preservation of animals or plants; as animal \textit{economy}; vegetable \textit{economy}.
						\item[7] Distribution or due order of things.
						\item[8] Judicious and frugal management of public affairs; as political \textit{economy}.
						\item[9] System of management; general regulation and disposition of the affairs of a state or nation, or of any department of government.
					\end{compactitem}
				All of these senses could be in play here, but the most important ones are probably \#5, \#6, and \#9. Note that all three of these allow us to connect the idea of an ``economy'' here to BY's understanding of the priesthood as the system of laws that govern everything.
				
				Calling this a vision on ``the economy of God'' tells us something important about how we are supposed to take it; it isn't just about the kingdoms. There is much more here.
				
				The only place that ``economy'' appears in the scriptures is \hyperref[DC77-6]{DC 77:6}.
				}%
		HEAR, O ye heavens, and give ear, O earth,%
			\footnote{%
				\label{DC76-1-hear}%
				This invocation at the beginning, which calls out to the heavens and to the earth, comes from \hyperref[ISAIAH-1-2]{Isaiah 1:2}, the beginning of Isaiah's vision: ``Hear, O heavens, and give ear, O earth: for the LORD hath spoken, I have nourished and brought up children, and they have rebelled against me.'' The ``hear...give ear'' construction appears elsewhere in the OT, including again in Isaiah 1 (\hyperref[ISAIAH-1-10]{verse 10}) and other places in Isaiah (\hyperref[ISAIAH-28-23]{28:23}, \hyperref[ISAIAH-32-9]{32:9}, \hyperref[ISAIAH-42-23]{42:23}, and sort of \hyperref[ISAIAH-51-4]{51:4}), and also in \hyperref[DEUTERONOMY-32-1]{Deuteronomy 32:1}, \hyperref[JUDGES-5-3]{Judges 5:3}, \hyperref[JOB-34-2]{Job 34:2}, plus in \hyperref[PSALMS]{Psalms}, \hyperref[JEREMIAH]{Jeremiah}, \hyperref[HOSEA]{Hosea}, and \hyperref[JOEL]{Joel}.
				
				They key here is that the speaker---who appears to be a combination of JS and Sidney Rigdon, who experienced the vision---is not the Lord, but addresses the heavens and earth anyway. See \hyperref[ISAIAH-1-2]{Isaiah 1:2} for more discussion. Note that in verse 1 of that chapter in Isaiah, the text introduces itself as ``the vision.''
				
				The occurrence in \hyperref[DEUTERONOMY-32-1]{Deuteronomy 32:1} is particularly interesting, as that is the ``Song of Moses'' and the very next verse is referenced by another important (but much later) passage in the DC, \hyperref[DC121-45]{DC 121:45}.
				}
		and rejoice ye inhabitants thereof, for the Lord is God, and beside%
			\footnote{%
				\label{DC76-1-beside}%
				See \hyperref[DEUTERONOMY-4-35]{Deuteronomy 4:35}, ``Unto thee it was shewed, that thou mightest know that the LORD he is God; there is none else beside him,'' and \hyperref[ISAIAH-45-21]{Isaiah 45:21}, ``Tell ye, and bring them near; yea, let them take counsel together: who hath declared this from ancient time? who hath told it from that time? have not I the LORD? and there is no God else beside me; a just God and a Saviour; there is none beside me''; see also \hyperref[1SAMUEL-2-2]{1 Samuel 2:2}, \hyperref[2SAMUEL-7-22]{2 Samuel 7:22}, \hyperref[1KINGS-8-60]{1 Kings 8:60}, \hyperref[1CHRONICLES-17-20]{1 Chronicles 17:20}, \hyperref[ISAIAH-43-11]{Isaiah 43:11}, \hyperref[ISAIAH-44-6]{Isaiah 44:6}, \hyperref[ISAIAH-45-5]{Isaiah 45:5}, \hyperref[HOSEA-13-4]{Hosea 13:4}, and \hyperref[MOSES-1-6]{Moses 1:6}.
				
				The ``beside him'' construction is pretty common in the OT, as you can see from these references; the word ``besides'' in any form appears only three times in the BOM, four times in the DC, and once in the PGP.
				
				When you combine this construction with the first half of the verse, it's a very biblical way of calling on the hearers to pay attention, because something very important is about to come out.
				}
		him there is no Savior.%
			\footnote{%
				\label{DC76-1-savior}%
				Note that the JSP version as well as the original EMS printing have this as ``and beside him there is none else'' (the comparison after the verse). The ``no Savior'' bit gets added in the 1835 edition of the DC.
				}

		\begin{framed}
		\scriptsize
			\begin{compactdesc}
				\item[JSP] Here O ye heavens \& give ere [ear] O earth and rejoice ye inhabitants thereof for the lord he is God and beside him there is none else
				\item[EMS] HEAR, O ye Heavens, and give ear, O earth, and rejoice ye inhabitants thereof, for the Lord he is God, and beside him there is none else;
				\item[BC] ---
				\item[DC 1835] Hear, O ye heavens, and give ear, O earth, and rejoice ye inhabitants thereof, for the Lord is God, and beside him there is no Savior;
				\item[DC 1844] Hear, O ye heavens, and give ear, O earth, and rejoice ye inhabitants thereof, for the Lord is God, and beside him there is no Savior;
			\end{compactdesc}
		\end{framed}

	% -----
	\subsubsection{Section 76:2} % *DC76-2 
	% -----
		\label{DC76-2}

		Great is his wisdom, marvelous are his ways,%
			\footnote{%
				\label{DC76-2-marvelous}%
				Verse 1 includes a lot of biblical language, but then there's an immediate shift here---the ``great/marvelous'' combination occurs only twice in the OT, at \hyperref[JOB-5-9]{Job 5:9} and \hyperref[JOB-37-5]{Job 37:5} (sort of; note that the KJV spells the word as ``marvellous''). You then get a similar construction in \hyperref[REVELATIONNT-15-1]{Revelation 15:1} and \hyperref[REVELATIONNT-15-3]{15:3} (verse 3 again mentions the ``song of Moses''). But then the combination of these words appears dozens of times in the BOM. So the first verse has a very biblical construction, but then in this second verse, or second part (remember the verses came much later), you get a very restoration-y construction. Note that ``extent'' later in this verse only appears twice in the scriptures, here and \hyperref[DC19-38]{DC 19:38}.
				}
		and the extent of his doings none can find out.

		% \begin{framed}
		% \scriptsize
			% \begin{compactdesc}
				% \item[JSP] 
				% % \item[BC] 
				% % \item[]
			% \end{compactdesc}
		% \end{framed}

	% -----
	\subsubsection{Section 76:3} % *DC76-3 
	% -----
		\label{DC76-3}

		His purposes fail not, neither are there any who can stay his hand.%
			\footnote{%
				\label{DC76-3-stay}%
				The ``stay...hand'' construction appears a couple of times in the OT (see especially \hyperref[DANIEL-4-35]{Daniel 4:35}, but also \hyperref[2SAMUEL-24-16]{2 Samuel 24:16} and \hyperref[1CHRONICLES-21-15]{1 Chronicles 21:15}). It appears a few times in the BOM (\hyperref[ALMA-10-23]{Alma 10:23} and \hyperref[MORMON-8-26]{Mormon 8:26} are two examples) and in the DC too (see \hyperref[DC38-22]{DC 38:22 and 33}).
				}

		% \begin{framed}
		% \scriptsize
			% \begin{compactdesc}
				% \item[JSP] 
				% % \item[BC] 
				% % \item[]
			% \end{compactdesc}
		% \end{framed}

	% -----
	\subsubsection{Section 76:4} % *DC76-4 
	% -----
		\label{DC76-4}

		From eternity to eternity he is the same, and his years never fail.%
			\footnote{%
				\label{DC76-4-fail}%
				See \hyperref[HEBREWS-1-12]{Hebrews 1:12}: ``And as a vesture shalt thou fold them up, and they shall be changed: but thou art the same, and thy years shall not fail.''
				}

		% \begin{framed}
		% \scriptsize
			% \begin{compactdesc}
				% \item[JSP] 
				% % \item[BC] 
				% % \item[]
			% \end{compactdesc}
		% \end{framed}

	% -----
	\subsubsection{Section 76:5} % *DC76-5 
	% -----
		\label{DC76-5}

		For thus saith the Lord---I, the Lord, am merciful and gracious%
			\footnote{%
				\label{DC76-5-gracious}%
				Here, in the context of the Lord having these qualities ``unto those that fear'' him, I think we want to take ``merciful'' as ``mercy-giving'' and ``gracious'' as ``grace-giving.'' See \hyperref[NEHEMIAH-9-17]{Nehemiah 9:17} for the pair of ``merciful and gracious.''
				}
		unto those who fear me,%
			\footnote{%
				% \label{}%
				Compare \hyperref[LUKE-1-50]{Luke 1:50}.
				}
		and delight to honor%
			\footnote{%
				\label{DC76-5-delight}%
				The ``delight...honor'' construction appears only here, but see \hyperref[ISAIAH-42-1]{Isaiah 42:1}.
				}
		those who serve me in righteousness and in truth unto the end.

		% \begin{framed}
		% \scriptsize
			% \begin{compactdesc}
				% \item[JSP] 
				% % \item[BC] 
				% % \item[]
			% \end{compactdesc}
		% \end{framed}

	% -----
	\subsubsection{Section 76:6} % *DC76-6 
	% -----
		\label{DC76-6}

		Great shall be their reward%
			\footnote{%
				% \label{}%
				Augustine, DDC I:32: ``For what words can express, what thoughts conceive, the reward which he is going to give at the end, seeing that he has already given us, to support us on our journey, so much of his spirit, in order that in the troubles of this life we may have this enormous confidence and delight in one whom we do not yet behold, and seeing that he has also bestowed individual gifts for the consolidation of his church, in order that we may perform the tasks that he has indicated not only without murmuring but even with positive enjoyment?''
				}
		and eternal shall be their glory.

		% \begin{framed}
		% \scriptsize
			% \begin{compactdesc}
				% \item[JSP] 
				% % \item[BC] 
				% % \item[]
			% \end{compactdesc}
		% \end{framed}

	% -----
	\subsubsection{Section 76:7} % *DC76-7 
	% -----
		\label{DC76-7}

		And to them%
			\footnote{%
				\label{DC76-7-them}%
				Remember, the ``them'' here goes back to verse 5: we're talking about ``those who fear'' the Lord, and ``those who serve [him] in righteousness and truth unto the end.''
				}
		will I reveal all \hyperref[MYSTERY]{mysteries},%
			\footnote{%
				\label{DC76-7-mysteries}%
				This verse and the next three are key passages for the \hyperref[EPIST-REVELATION]{epistemology of revelation}. They have a lot in common with other passages of scripture, as I've noted in various places, but there is a particularly strong parallel with \hyperref[1CORINTHIANS-2]{1 Corinthians 2}.
				
				The word ``mystery'' doesn't appear in the OT; it appears a bunch of times in the \hyperref[MYSTERY-NT]{NT}. It doesn't appear there, though, in conjunction with the word ``reveal''; you only get that in the restoration scriptures, and there only \hyperref[JACOB-4-8]{Jacob 4:8}, \hyperref[ALMA-26-22]{Alma 26:22}, this verse, and also \hyperref[DC77-6]{DC 77:6}.
				}
		yea, all the hidden \hyperref[MYSTERY]{mysteries} of my kingdom from days of old, and for ages to come, will I make known unto them the good pleasure of my will%
			\footnote{%
				% \label{}%
				See \hyperref[EPHESIANS-1-9]{Ephesians 1:9} and \hyperref[DC38-39]{DC 38:39} plus notes there.
				}
		concerning all things%
			\footnote{%
				\label{DC76-7-will}%
				This is a weird part of the verse---it feels like it's lacking punctuation or something to be clearer about what is being said. It looks like the direct object of ``will I make known unto them'' should be ``all the hidden mysteries,'' but then right after it you get ``the good pleasure of my will,'' which seems like it ought to be the direct object. So I'm not totally sure what is going on here. It could be punctuated in different ways to make it make sense, but switching the order of ``I'' and ``will'' make it a bit difficult:
					\begin{compactitem}
						\item[\#1] And to them will I reveal all mysteries, yea, all the hidden mysteries of my kingdom from days of old, and for ages to come. Will I make known unto them the good pleasure of my will concerning all things pertaining to my kingdom.
						\item[\#2] And to them will I reveal all mysteries, yea, all the hidden mysteries of my kingdom from days of old, and for ages to come, will I make known unto them; the good pleasure of my will concerning all things pertaining to my kingdom.
					\end{compactitem}
				I think I prefer \#2, because it allows you to put it together better with the next verse as well; the verse break is probably not a good one. So the two verses together would read like this:
					\begin{quote}
						And to them will I reveal all mysteries, yea, all the hidden mysteries of my kingdom from days of old, and for ages to come, will I make known unto them. The good pleasure of my will concerning all things pertaining to my kingdom, yea, even the wonders of eternity shall they know, and things to come will I show them, even the things of many generations.
					\end{quote}
				This reading seems to make the most sense, and also allows you to use the ``yea'' properly, which introduces either an explanatory apposition or just an additional (``and also'') thing that will be revealed.
				}
		pertaining to my kingdom.%
			\footnote{%
				\label{DC76-7-kingdom}%
				Note the difference in the JSP version, which finishes, ``concerning all things to come.'' 
				}

		\begin{framed}
		\scriptsize
			\begin{compactdesc}
				\item[JSP] \textbf{and unto them will I reveal all my misteries yea all the hiden misteries of my Kingdom from days of old and for ages to come will I make Known unto them the good pleasure of my will concerning all things to come}
				% \item[BC] 
				% \item[]
			\end{compactdesc}
		\end{framed}

	% -----
	\subsubsection{Section 76:8} % *DC76-8 
	% -----
		\label{DC76-8}

		Yea, even the wonders of eternity shall they know, and things to come will I show them, even the things of many generations.

		\begin{framed}
		\scriptsize
			\begin{compactdesc}
				\item[JSP] \textbf{yea} even the wonders of eternity shall they know and things to come will I shew them even the things of many generations
				% \item[BC] 
				% \item[]
			\end{compactdesc}
		\end{framed}

	% -----
	\subsubsection{Section 76:9} % *DC76-9 
	% -----
		\label{DC76-9}

		And their wisdom shall be great, and their understanding reach to heaven; and before them the wisdom of the wise shall perish, and the understanding of the prudent shall come to naught.%
			\footnote{%
				\label{DC76-9-prudent}%
				This verse is quoting \hyperref[ISAIAH-29-14]{Isaiah 29:14}, though it makes a slight change---the end of that verse reads, ``the wisdom of their wise men shall perish, and the understanding of their prudent men shall be hid,'' but here it says that the understanding of the prudent ``shall come to naught.'' See also \hyperref[ISAIAH-5-20]{Isaiah 5:20--21}, \hyperref[HOSEA-14-9]{Hosea 14:9}, \hyperref[MATTHEW-11-25]{Matthew 11:25} and \hyperref[LUKE-10-21]{Luke 10:21}, \hyperref[1CORINTHIANS-1-19]{1 Corinthians 1:19}, and \hyperref[2NEPHI-9-43]{2 Nephi 9:43}.
				}

		\begin{framed}
		\scriptsize
			\begin{compactdesc}
				\item[JSP] there wisdom shall be great and there understanding reach to heaven and before them the wisdom of the wise shall perish and the understanding of the prudent shall come to naught
				% \item[BC] 
				% \item[]
			\end{compactdesc}
		\end{framed}

	% -----
	\subsubsection{Section 76:10} % *DC76-10 
	% -----
		\label{DC76-10}

		For by my Spirit will I enlighten them, and by my power will I make known unto them the secrets of my will---yea, even those things which eye has not seen, nor ear heard, nor yet entered into the heart of man.%
			\footnote{%
				\label{DC76-10-heart}%
				See \hyperref[1CORINTHIANS-2]{1 Corinthians 2}, especially verse \hyperref[1CORINTHIANS-2-9]{9}. The poetic version, in stanza 10, also makes the connection to the ``natural man'' as discussed in that chapter of Corinthians.
				} 

		\begin{framed}
		\scriptsize
			\begin{compactdesc}
				\item[JSP] for by my spirit will I enlighten them and by my power will I make known unto them the secrets of my will yea even those things which eye has not seen nor ear heard nor yet entered into the heart of man
				% \item[BC] 
				% \item[]
			\end{compactdesc}
		\end{framed}

	% -----
	\subsubsection{Section 76:11} % *DC76-11 
	% -----
		\label{DC76-11}

		We, Joseph Smith, Jun., and Sidney Rigdon, being in the Spirit on the sixteenth day of February, in the year of our Lord one thousand eight hundred and thirty-two---

		% \begin{framed}
		% \scriptsize
			% \begin{compactdesc}
				% \item[JSP] 
				% % \item[BC] 
				% % \item[]
			% \end{compactdesc}
		% \end{framed}

	% -----
	\subsubsection{Section 76:12} % *DC76-12 
	% -----
		\label{DC76-12}

		By the power of the Spirit%
			\footnote{%
				% \label{}%
				For ``power of the Spirit,'' see \hyperref[POWER-PHRASES]{here}.
				}
		our eyes were opened and our understandings were enlightened, so as to see and understand the things of God---%
			\footnote{%
				\label{DC76-12-spirit}%
				This verse then describes how JS and Rigdon experienced what the Lord himself was describing earlier---being enlightened by the Spirit and seeing and understanding the things of God. They are living the promise that the Lord gave earlier.
				}

		% \begin{framed}
		% \scriptsize
			% \begin{compactdesc}
				% \item[JSP] 
				% % \item[BC] 
				% % \item[]
			% \end{compactdesc}
		% \end{framed}

	% -----
	\subsubsection{Section 76:13} % *DC76-13 
	% -----
		\label{DC76-13}

		Even those things which were from the beginning before the world was, which were \hyperref[ORDINATION]{ordained} of the Father, through his Only Begotten Son,%
			\footnote{%
				\label{DC76-13-son}%
				The Son being the agent of the Father, to come here and do his will; see \hyperref[JOHN-4-34]{John 4:34} and \hyperref[JOHN-5-30]{John 5:30}. ``Ordain'' here could be meaning \#2, \#3, or \#5 from the 1828 dictionary---\#3 seems most appropriate, given the context of the things existing for such a long time prior to the present moment.
				}
		who was in the bosom%
			\footnote{%
				\label{DC76-13-bosom}%
				``Bosom'' occurs here and also in verses \hyperref[DC76-25]{25} and \hyperref[DC76-39]{39}, at least in this section. The key verse in the NT is \hyperref[JOHN-1-18]{John 1:18}; see the notes there. The meaning in a passage like this one is intended to convey the very closest relationship between these two individuals.
				}
		of the Father, even from the beginning;

		% \begin{framed}
		% \scriptsize
			% \begin{compactdesc}
				% \item[JSP] 
				% % \item[BC] 
				% % \item[]
			% \end{compactdesc}
		% \end{framed}

	% -----
	\subsubsection{Section 76:14} % *DC76-14 
	% -----
		\label{DC76-14}

		Of whom we bear record; and the record which we bear is the fulness of the gospel of Jesus Christ,%
			\footnote{%
				\label{DC76-14-gospel}%
				This is one of the scriptural ``definitions'' of the term ``gospel''; see \hyperref[GOSPEL-definition]{this study} for more.
				}
		who is the Son, whom we saw and with whom we conversed in the heavenly vision.

		% \begin{framed}
		% \scriptsize
			% \begin{compactdesc}
				% \item[JSP] 
				% % \item[BC] 
				% % \item[]
			% \end{compactdesc}
		% \end{framed}

	% -----
	\subsubsection{Section 76:15} % *DC76-15 
	% -----
		\label{DC76-15}

		For while we were doing the work of translation,%
			\footnote{%
				\label{DC76-15-translation}%
				The next few verses are helpful for seeing how JS himself understood the ``translation'' of the Bible. See \hyperref[TRANSLATION-meaning]{this study} at \hyperref[TRANSLATION]{Translation} for more sources related to this issue.
				}
		which the Lord had appointed unto us, we came to the twenty-ninth verse of the fifth chapter of John, which was given unto us as follows---

		% \begin{framed}
		% \scriptsize
			% \begin{compactdesc}
				% \item[JSP] 
				% % \item[BC] 
				% % \item[]
			% \end{compactdesc}
		% \end{framed}

	% -----
	\subsubsection{Section 76:16} % *DC76-16 
	% -----
		\label{DC76-16}

		Speaking of the resurrection of the dead, concerning those who shall hear the voice of the Son of Man:

		% \begin{framed}
		% \scriptsize
			% \begin{compactdesc}
				% \item[JSP] 
				% % \item[BC] 
				% % \item[]
			% \end{compactdesc}
		% \end{framed}

	% -----
	\subsubsection{Section 76:17} % *DC76-17 
	% -----
		\label{DC76-17}

		And shall come forth; they who have done good in the resurrection of the just; and they who have done evil, in the resurrection of the unjust.%
			\footnote{%
				\label{DC76-17-john5}%
				The KJV gives \hyperref[JOHN-5-29]{John 5:29} as: ``And shall come forth; they that have done good, unto the resurrection of life; and they that have done evil, unto the resurrection of damnation.'' Here is the Greek: \gr{καὶ ἐκπορεύσονται οἱ τὰ ἀγαθὰ ποιήσαντες εἰς ἀνάστασιν ζωῆς, οἱ δὲ τὰ φαῦλα πράξαντες εἰς ἀνάστασιν κρίσεως.} Note that \hyperref[LUKE-14-14]{Luke 14:14} speaks of the ``resurrection of the just'' (\gr{ἐν τῇ ἀναστάσει τῶν δικαίων}), and so does \hyperref[ACTS-24-15]{Acts 24:15}. The ``just'' are those who are \hyperref[DIKAIOS]{\grl{δίκαιος}}, who are ``just'' and especially ``lawful.'' See also \hyperref[ALMA-12-8]{Alma 12:8}; \hyperref[3NEPHI-26-5]{3 Nephi 26:5} (Christ speaking) speaks of the resurrection as one of ``everlasting life'' and one of ``damnation,'' more like the KJV John 5:29. Maybe the best approach to take is to just put them together: those who are \gr{δίκαιος} get a resurrection of everlasting life, and those who are not \gr{ἄδικος} get a resurrection of damnation. This would be in line with BY's view of how the law of the priesthood is intended to preserve life. The ``resurrection of the unjust'' is still a resurrection, though, so life is still being preserved; you're obeying some of the laws, and not trying to be a law unto yourself, maybe. 
				}

		% \begin{framed}
		% \scriptsize
			% \begin{compactdesc}
				% \item[JSP] 
				% % \item[BC] 
				% % \item[]
			% \end{compactdesc}
		% \end{framed}

	% -----
	\subsubsection{Section 76:18} % *DC76-18 
	% -----
		\label{DC76-18}

		Now this caused us to marvel, for it was given unto us of the Spirit.

		% \begin{framed}
		% \scriptsize
			% \begin{compactdesc}
				% \item[JSP] 
				% % \item[BC] 
				% % \item[]
			% \end{compactdesc}
		% \end{framed}

	% -----
	\subsubsection{Section 76:19} % *DC76-19 
	% -----
		\label{DC76-19}

		And while we meditated upon these things, the Lord touched the eyes of our understandings and they were opened,%
		\footnote{%
			\label{DC76-19-eyes}%
			The ``eyes of the understanding'' locution comes from \hyperref[EPHESIANS-1-18]{Ephesians 1:18}, but both the Greek and modern translations differ from the way that the KJV puts it. The locution also appears elsewhere in the DC, including at \hyperref[DC110-1]{110:1} and \hyperref[DC138-11]{138:11}. This verse is also a callback to \hyperref[DC76-12]{verse 12} from before.
			} 
		and the glory of the Lord shone round about.

		% \begin{framed}
		% \scriptsize
			% \begin{compactdesc}
				% \item[JSP] 
				% % \item[BC] 
				% % \item[]
			% \end{compactdesc}
		% \end{framed}

	% -----
	\subsubsection{Section 76:20} % *DC76-20 
	% -----
		\label{DC76-20}

		And we beheld the glory of the Son, on the right hand of the Father, and received of his \hyperref[FULNESS]{fulness};%
			\footnote{%
				\label{DC76-20-fulness}%
				There are many verses relevant to understanding what it means to ``receive of his fulness'' here, including \hyperref[JOHN-1-16]{John 1:16} (\gr{ἐκ τοῦ} \hyperref[PLEROMA]{\grl{πληρώματος}} \gr{αὐτοῦ ἡμεῖς πάντες ἐλάβομεν, καὶ χάριν ἀντὶ χάριτος}), \hyperref[DC63-20]{DC 63:20--21}, \hyperref[DC76-56]{76:56}, \hyperref[DC76-71]{76:71--94}, especially verses \hyperref[DC76-71]{71--77}; \hyperref[DC88-29]{88:29--30}; \hyperref[DC93-12]{93:12--35}, and \hyperref[DC132-6]{132:6}. See also the NT verses listed in the LSJ entry for \hyperref[PLEROMA]{\grl{πλήρωμα}}.
				}

		% \begin{framed}
		% \scriptsize
			% \begin{compactdesc}
				% \item[JSP] 
				% % \item[BC] 
				% % \item[]
			% \end{compactdesc}
		% \end{framed}

	% -----
	\subsubsection{Section 76:21} % *DC76-21 
	% -----
		\label{DC76-21}

		And saw the holy angels, and them who are \hyperref[SANCTIFICATION]{sanctified} before his throne, worshiping God, and the Lamb, who worship him forever and ever.

		% \begin{framed}
		% \scriptsize
			% \begin{compactdesc}
				% \item[JSP] 
				% % \item[BC] 
				% % \item[]
			% \end{compactdesc}
		% \end{framed}

	% -----
	\subsubsection{Section 76:22} % *DC76-22 
	% -----
		\label{DC76-22}

		And now, after the many \hyperref[TESTIMONY]{testimonies} which have been given of him, this is the \hyperref[TESTIMONY]{testimony}, last of all, which we give of him: That he lives!

		% \begin{framed}
		% \scriptsize
			% \begin{compactdesc}
				% \item[JSP] 
				% % \item[BC] 
				% % \item[]
			% \end{compactdesc}
		% \end{framed}

	% -----
	\subsubsection{Section 76:23} % *DC76-23 
	% -----
		\label{DC76-23}

		For we saw him, even on the right hand of God; and we heard the voice bearing record that he is the Only Begotten of the Father---

		% \begin{framed}
		% \scriptsize
			% \begin{compactdesc}
				% \item[JSP] 
				% % \item[BC] 
				% % \item[]
			% \end{compactdesc}
		% \end{framed}

	% -----
	\subsubsection{Section 76:24} % *DC76-24 
	% -----
		\label{DC76-24}

		That by him, and through him, and of him,%
			\footnote{%
				\label{DC76-24-by}%
				See \hyperref[ROMANS-11-36]{Romans 11:36}, \hyperref[DC88-41]{DC 88:41}, \hyperref[DC93-10]{93:10}.
				}
		the worlds are and were created, and the inhabitants thereof are begotten sons and daughters unto God.%
			\footnote{%
				\label{DC76-24-unto}%
				This verse could be the basis for believing that we are not necessarily of the same form as God---we are begotten sons and daughters \textit{unto} him, not \textit{of} him. This is important because it's a way to think that God doesn't necessarily have human form. Of course, other scriptures and most of LDS theology pushes against this, but it's very important for thinking about biological evolution. See also this unusual comment in the opening part of the \hyperref[TEMPLE-ENDOWMENT-ENDOW-PRE-1990-INTRO]{endowment} prior to 1990: ``The work of the six creative periods will be represented. They will organize man in their own likeness and image, male and female. This, however, is simply figurative, so far as the man and the woman are concerned.'' That last sentence was removed in the 1990 revision.
				
				See also \hyperref[DC88-20]{DC 88:20}, with its use of ``bodies'';
				
				See also Orson Pratt \hyperref[EMS-1852-OP-28-AUG-1852]{here}, 17a, where he discusses this verse and the ``unto.''
				}

		% \begin{framed}
		% \scriptsize
			% \begin{compactdesc}
				% \item[JSP] 
				% % \item[BC] 
				% % \item[]
			% \end{compactdesc}
		% \end{framed}

	% -----
	\subsubsection{Section 76:25} % *DC76-25 
	% -----
		\label{DC76-25}

		And this we saw also, and bear record, that an angel of God who was in authority%
			\footnote{%
				\label{DC76-25-authority}%
				If this revelation had been received in the Greek of the NT, the word here would be \hyperref[EXOUSIA]{\grl{ἐξουσία}} for ``authority.'' See the notes at \hyperref[DC76-31]{verse 31} below.
				}
		in the presence of God, who rebelled against the Only Begotten Son whom the Father loved and who was in the \hyperref[BOSOM]{bosom}%
			\footnote{%
				% \label{}%
				See notes at \hyperref[JOHN-1-18]{John 1:18}.
				}
		of the Father, was thrust down from the presence of God and the Son,

		\begin{framed}
		\scriptsize
			\begin{compactdesc}
				\item[JSP] and this we saw also and bear reccord that an angel of God who was in authority in the presence of God who rebelled against the only begotten son whom the father loved who was in the bosom with the father and was thrust down from the presence of \sout{the father} God and the son
				% \item[BC] 
				% \item[]
			\end{compactdesc}
		\end{framed}

	% -----
	\subsubsection{Section 76:26} % *DC76-26 
	% -----
		\label{DC76-26}

		And was called \hyperref[SONSOFPERDITION]{Perdition},%
			\footnote{%
				% \label{}%
				See \hyperref[JOHN-17-12]{John 17:12} but especially \hyperref[2THESSALONIANS-2-3]{2 Thessalonians 2:3}---the end of that verse says, ``and that man of sin be revealed, the son of perdition.'' The ``man of sin'' is \gr{ὁ ἄνθρωπος τῆς} \hyperref[ANOMIA]{\grl{ἀνομίας}} (there is a different manuscript tradition that does have \gr{ἁμαρτίας}; the reading with \gr{ἀνομίας} has a certain level of ``B,'' ``almost certain''). This is extremely important because it solidifies the connection to \hyperref[DC88-35]{DC 88:35} as a description of Satan, and as an explanation of why he could not have succeeded: ``That which breaketh a law, and abideth not by law, but seeketh to become a law unto itself, and willeth to abide in sin, and altogether abideth in sin, cannot be sanctified by law, neither by mercy, justice, nor judgment. Therefore, they must remain filthy still.'' See also \hyperref[DC76-33-]{verse 33} below.
				}
		for the heavens wept over him---he was \hyperref[LUCIFER]{Lucifer},%
			\footnote{%
				% \label{}%
				See the important discussion of the title ``Lucifer'' at \hyperref[ISAIAH-14-12]{Isaiah 14:12}. The short version is that modern interpretations and translations seem to agree that those verses in Isaiah are not referring to the fall of Satan, and are referring to human pride or to a Babylonian king (as possibly a symbol of human pride). JS clearly took it to refer to Satan, however, since he quotes that passage here in discussing Satan. 
				
				The methodology here is a bit tenuous here, because we have go to with JS in taking Isaiah to be talking about Satan, but we can look at the Hebrew for a bit of support for the theology that JS builds up around Satan here and in other places. In the previous verse we read that Satan was an ``angel of God'' who was ``in authority in the presence of God.'' In the Isaiah verse, the term which gets translated as ``Lucifer'' in the KJV is \he{heylel}, coming from a root that means ``to shine'' (see discussion in my notes at that verse). If you connect to restoration ideas like \hyperref[DC50-24]{DC 50:24} and \hyperref[DC93-28]{93:28}, then you could take Satan's ``shining'' as an example of how he had gained a lot of light and truth but for some reason chose to go away from that.
				}
		a son of the morning.

		\begin{framed}
		\scriptsize
			\begin{compactdesc}
				\item[JSP] and was called perdition for the heavens wept over him for he was Lucifer even the son of the morning
				% \item[BC] 
				% \item[]
			\end{compactdesc}
		\end{framed}

	% -----
	\subsubsection{Section 76:27} % *DC76-27 
	% -----
		\label{DC76-27}

		And we beheld, and lo, he is fallen! is fallen, even a son of the morning!

		% \begin{framed}
		% \scriptsize
			% \begin{compactdesc}
				% \item[JSP] 
				% % \item[BC] 
				% % \item[]
			% \end{compactdesc}
		% \end{framed}

	% -----
	\subsubsection{Section 76:28} % *DC76-28 
	% -----
		\label{DC76-28}

		And while we were yet in the Spirit, the Lord commanded us that we should write the vision; for we beheld \hyperref[SATAN]{Satan}, that old serpent, even the devil, who rebelled against God,%
			\footnote{%
				% \label{}%
				See \hyperref[MOSES-4-3]{Moses 4:3--4}, where God is speaking: ``Satan rebelled against me.''
				}
		and sought to take the kingdom%
			\footnote{%
				% \label{}%
				The concept of a ``kingdom'' is closely connected to ideas of laws in the restoration, so you could read this as saying that Satan was trying to change or go against the laws that form the basis of the priesthood (in BY's way of understanding the priesthood); see \hyperref[DC88-35]{DC 88:35} again, but also \hyperref[DC88-36]{36--38}.
				}
		of our God and his Christ---

		% \begin{framed}
		% \scriptsize
			% \begin{compactdesc}
				% \item[JSP] 
				% % \item[BC] 
				% % \item[]
			% \end{compactdesc}
		% \end{framed}

	% -----
	\subsubsection{Section 76:29} % *DC76-29 
	% -----
		\label{DC76-29}

		Wherefore, he maketh war with the saints of God, and encompasseth
		them round about.%
			\footnote{%
				\label{DC76-29-encompasseth}%
				See note at \hyperref[2NEPHI-4-18]{2 Nephi 4:18}.
				}

		% \begin{framed}
		% \scriptsize
			% \begin{compactdesc}
				% \item[JSP] 
				% % \item[BC] 
				% % \item[]
			% \end{compactdesc}
		% \end{framed}

	% -----
	\subsubsection{Section 76:30} % *DC76-30 
	% -----
		\label{DC76-30}

		And we saw a vision of the sufferings of those with whom he made war and overcame,%
			\footnote{%
				% \label{}%
				The battle is ongoing---this is a vision involving people who lost the war in the long run. The rest of us will fight this war and lose some battles but not be overcome in the end; see \hyperref[HELAMAN-5-12]{Helaman 5:12}, and in connection with that verse, \hyperref[2NEPHI-1-13]{2 Nephi 1:13}, \hyperref[2NEPHI-26-22]{2 Nephi 26:22}, \hyperref[ALMA-12-11]{Alma 12:11}, \hyperref[ALMA-34-35]{Alma 34:35} and \hyperref[ALMA-34-39]{39}, and \hyperref[HELAMAN-3-29]{Helaman 3:29}. We \textit{allow} this to happen to us, by the ways we live and think over long periods of time (see the next verse, and also \hyperref[ALMA-5-20]{Alma 5:20} and \hyperref[MORONI-7-17]{Moroni 7:17}).
				
				See also \hyperref[DC50-8]{DC 50:8} and notes there, 
				}
		for thus came the voice of the Lord unto us:

		\begin{framed}
		\scriptsize
			\begin{compactdesc}
				\item[JSP] and we saw a vision of the eternal sufferings of those with whom he made war and overcame for thus came the voice of the Lord unto us
				% \item[BC] 
				% \item[]
			\end{compactdesc}
		\end{framed}

	% -----
	\subsubsection{Section 76:31} % *DC76-31 
	% -----
		\label{DC76-31}

		Thus saith the Lord concerning all those who know my power, and have been made partakers thereof,%
			\footnote{%
				% \label{}%
				We've got three uses of ``power'' in this verse, and when I see this word in the restoration scriptures, I think to myself, what would the Greek word be if this verse appeared in the NT? ``Partake'' and ``power'' don't appear often together. The clearest example of their use together, besides this verse, is \hyperref[1CORINTHIANS-9-12]{1 Corinthians 9:12}, and the word there is \gr{ἐξουσία}, but it isn't really talking about the same idea, as far as I can tell. But when you look at the next bit in this verse---about people ``suffering themselves'' (allowing themselves) through the devil's power to be overcome---then it looks more like \gr{ἐξουσία} is the right word, rather than \gr{δύναμις} (note that \hyperref[HEBREWS-2-14]{Hebrews 2:14} mentions a ``power of death,'' where the ``power'' is \gr{κράτος}). Then you get to the end of the verse and see that these people deny the truth and defy God's power, which definitely sounds like \gr{ἐξουσία} in connection with \hyperref[DC88-35]{DC 88:35}, but could also be \gr{δύναμις} too (compare \hyperref[MOSES-4-3]{Moses 4:3}). Note the connection between \hyperref[DC76-34]{DC 76:34}, \hyperref[MATTHEW-12-32]{Matthew 12:32}, and \hyperref[DC84-41]{DC 84:41}; these verses help us understand that Satan must have been a partaker in the priesthood (the \gr{ἐξουσία}, through which he gained the \gr{δύναμις}). He then broke the covenant and so lost the power.
				} 
		and suffered themselves through the power of the devil to be overcome, and to deny the truth%
			\footnote{%
				% \label{}%
				Compare JS in various accounts of his \hyperref[EMS-1844-JS-7-APR-1844]{7 April 1844} discourse---speaking of this kind of person, who deny the Lord in this special way, he says: ``he has got to say that the Sun does not shine while he sees it he has got to deny J. C. when the heavens are open to him'' (Thomas Bullock).
				}
		and defy my power---

		% \begin{framed}
		% \scriptsize
			% \begin{compactdesc}
				% \item[JSP] 
				% % \item[BC] 
				% % \item[]
			% \end{compactdesc}
		% \end{framed}

	% -----
	\subsubsection{Section 76:32} % *DC76-32 
	% -----
		\label{DC76-32}

		They are they who are the sons of perdition, of whom I say that it had been better for them never to have been born;%
		\footnote{%
			% \label{}%
			I take this to mean that it would have been better for these individuals to never come to Earth to receive a body. But why? Why would that have been better? I guess the question is about the ultimate difference in fate between that group of spirits who never did come to Earth or get a body, and those who did but then ended up in this group, of the ``sons of perdition.'' Now that I'm thinking about it, I guess I don't know anything about the fate of the former group.
			} 

		% \begin{framed}
		% \scriptsize
			% \begin{compactdesc}
				% \item[JSP] 
				% % \item[BC] 
				% % \item[]
			% \end{compactdesc}
		% \end{framed}

	% -----
	\subsubsection{Section 76:33} % *DC76-33 
	% -----
		\label{DC76-33}

		For they are vessels of wrath,%
			\footnote{%
				% \label{}%
				The phrase ``vessels of wrath'' comes from \hyperref[ROMANS-9-22]{Romans 9:22}. That verse is part of a passage speaking about different groups of people, or ``vessels'' (\gr{σκεύη}). There are some vessels which are vessels of wrath, and the way it is written allows us to read the verse as saying that they are filled with their \textit{own} wrath (which is how I've always taken this verse in DC) as well as with the wrath of God. The other group of people is the ``vessels of mercy,'' and along the same lines you can see them as filled with their own mercy (extended toward themselves and others) as well as with the mercy of God (\hyperref[DC88-40]{DC 88:40}, ``mercy hath compassion on mercy and claimeth her own''). The connection to the ``sons of perdition'' is made even clearer by the claim at the end of Romans 9:22, which says that these vessels are fitted ``to destruction,'' or \gr{εἰς ἀπώλειαν}. That's \hyperref[APOLEIA]{\grl{ἀπώλεια}}, the same word I was discussing up at \hyperref[DC76-26]{verse 26} above.
				}
		doomed to suffer the wrath of God, with the devil and his angels in eternity;

		% \begin{framed}
		% \scriptsize
			% \begin{compactdesc}
				% \item[JSP] 
				% % \item[BC] 
				% % \item[]
			% \end{compactdesc}
		% \end{framed}

	% -----
	\subsubsection{Section 76:34} % *DC76-34 
	% -----
		\label{DC76-34}

		Concerning whom I have said there is no forgiveness in this world nor in the world to come---%
			\footnote{%
				% \label{}%
				Presumably referring to \hyperref[MATTHEW-12-31]{Matthew 12:31--32} here, but see also \hyperref[DC42-18]{DC 42:18} and \hyperref[DC84-41]{DC 84:41}. I think that verse in DC 84 is especially important because it comes in context of a discussion on the priesthood, and then says that those who break the covenant ``shall not have forgiveness of sins in this world nor in the world to come,'' that is, the same language used to describe the individuals here in this verse in section 76.
				}

		% \begin{framed}
		% \scriptsize
			% \begin{compactdesc}
				% \item[JSP] 
				% % \item[BC] 
				% % \item[]
			% \end{compactdesc}
		% \end{framed}

	% -----
	\subsubsection{Section 76:35} % *DC76-35 
	% -----
		\label{DC76-35}

		Having denied the Holy Spirit after having received it,%
			\footnote{%
				% \label{}%
				A lot of scriptures talk about denying the Holy Ghost, including \hyperref[2NEPHI-28-4]{2 Nephi 28:4}, \hyperref[ALMA-39-5]{Alma 39:5--6}, \hyperref[MORONI-8-28]{Moroni 8:28}, and \hyperref[DC76-83]{DC 76:83}. There are also others that speak more specifically of denying the \textit{gift} of the Holy Ghost, which seems like a separate thing.
				}
		and having denied the Only Begotten Son of the Father,%
			\footnote{%
				% \label{}%
				Other verses to see about denying Christ include \hyperref[MATTHEW-10-33]{Matthew 10:33}, \hyperref[MATTHEW-26-34]{Matthew 26:34} (and the ensuing events for Peter), \hyperref[LUKE-12-9]{Luke 12:9}, \hyperref[ACTS-3-13]{Acts 3:13}, \hyperref[1JOHN-2-22]{1 John 2:22--23}, \hyperref[2NEPHI-25-28]{2 Nephi 25:28--29}, \hyperref[2NEPHI-28-32]{2 Nephi 28:32}, \hyperref[2NEPHI-31-14]{2 Nephi 31:14}, \hyperref[JACOB-7-9]{Jacob 7:9 and 19}, \hyperref[ALMA-30-39]{Alma 30:39}, \hyperref[3NEPHI-29-5]{3 Nephi 29:5}, \hyperref[4NEPHI-1-29]{4 Nephi 1:29}, \hyperref[MORMON-9-3]{Mormon 9:3}, \hyperref[MORONI-1-2]{Moroni 1:2--3}, \hyperref[MORONI-7-17]{Moroni 7:17}, \hyperref[MORONI-10-6]{Moroni 10:6}, \hyperref[DC76-43]{DC 76:43}, \hyperref[DC101-5]{DC 101:5}, and \hyperref[MOSES-6-28]{Moses 6:28}.
				}
		having crucified him unto themselves and put him to an open shame.%
			\footnote{%
				% \label{}%
				This last part of the verse comes from \hyperref[HEBREWS-6-6]{Hebrews 6:6}, which puts it this way: ``they crucify to themselves the Son of God afresh, and put him to an open shame.'' The verse in DC leaves out the part about ``afresh,'' or really ``again,'' which is supposed to come from the prefix on \gr{ἀνασταυρόω}, giving the idea of ``again.'' See the note at the Hebrews verse for more on that; the word technically just means ``crucify'' (without the ``again''), though the context seems to require the ``again,'' and it has long been understood that way.
				
				What does it mean, then, to crucify Christ unto yourself, and expose him to shame or ridicule? It could just be, you're refusing to acknowledge the atonement; you're refusing to let it be a part of your life, and maybe even, you're trying to cause it not to be a part of the lives of other people as well. You would have crucified Christ not necessariliy to hurt him, but to try to make his atonement of no effect in the world. I think that's the key here---it's not how cruel you would be to Christ as a person, it's what you would hope to accomplish by trying to erase Christ and his sacrifice from the world.
				}

		\begin{framed}
		\scriptsize
			\begin{compactdesc}
				\item[JSP] having denied the holy ghost after having received it and having denied the only begoten son of the fathe[r] crucifying him unto themselves and putting him to an open shame
				% \item[BC] 
				% \item[]
			\end{compactdesc}
		\end{framed}

	% -----
	\subsubsection{Section 76:36} % *DC76-36 
	% -----
		\label{DC76-36}

		These are they who shall go away into the lake of fire and brimstone, with the devil and his angels---

		% \begin{framed}
		% \scriptsize
			% \begin{compactdesc}
				% \item[JSP] 
				% % \item[BC] 
				% % \item[]
			% \end{compactdesc}
		% \end{framed}

	% -----
	\subsubsection{Section 76:37} % *DC76-37 
	% -----
		\label{DC76-37}

		And the only ones on whom the second death%
			\footnote{%
				% \label{}%
				For ``second death,'' see \hyperref[DEATH-PHRASES]{here}.
				}
		shall have any power;%
			\footnote{%
				% \label{}%
				The phrase ``second death'' occurs 12 times in the scriptures; it stems from the book of \hyperref[REVELATIONNT]{Revelation}, but is in the restoration scriptures too. See the complete list \hyperref[DEATH-PHRASES]{here}. You might see the ``second death'' as the ultimate fate of the sons of perdition, as described by BY, with the stuff about ``\hyperref[NATIVEELEMENT]{native element},'' but that doesn't really fit well with the way that the ``second death'' is spoken of elsewhere. See especially \hyperref[ALMA-12-15]{Alma 12:15--18} (and note the connection between that verse 18 and the next verse, 76:38) and \hyperref[REVELATIONNT-20-6]{Revelation 20:6--21:8}. The second death is to die ``as to things pertaining unto righteousness'' (Alma 12:16). There is a connection in other scriptures between righteousness and life (\hyperref[ROMANS-8-10]{Romans 8:10}, \hyperref[1NEPHI-15-36]{1 Nephi 15:36}, \hyperref[ALMA-5-58]{Alma 5:58}). Righteousness preserves life and organization, to put it in BY terms.
				} 

		% \begin{framed}
		% \scriptsize
			% \begin{compactdesc}
				% \item[JSP] 
				% % \item[BC] 
				% % \item[]
			% \end{compactdesc}
		% \end{framed}

	% -----
	\subsubsection{Section 76:38} % *DC76-38 
	% -----
		\label{DC76-38}

		Yea, verily, the only ones who shall not be redeemed in the due time of the Lord,%
			\footnote{%
				% \label{}%
				See \hyperref[MOSIAH-16-5]{Mosiah 16:5}, \hyperref[ALMA-11-41]{Alma 11:41}, \hyperref[ALMA-12-18]{Alma 12:18}, and \hyperref[MORONI-7-38]{Moroni 7:38}. The wicked will experience existence as though no redemption were made at all.
				}
		after the sufferings of his wrath.%
			\footnote{%
				% \label{}%
				Related to the idea of the ``buffetings of Satan''? See that phrase \hyperref[BUFFET-PHRASES]{here}.
				}

		\begin{framed}
		\scriptsize
			\begin{compactdesc}
				\item[JSP] yea verely the only ones who shall not be redeemed in the due time of the Lord after the sufferings of his wrath
				% \item[BC] 
				% \item[]
			\end{compactdesc}
		\end{framed}

	% -----
	\subsubsection{Section 76:39} % *DC76-39 
	% -----
		\label{DC76-39}

		For all the rest shall be brought forth by the resurrection of the dead, through the triumph and the glory of the Lamb, who was slain, who was in the bosom of the Father before the worlds were made.

		\begin{framed}
		\scriptsize
			\begin{compactdesc}
				\item[JSP] who shall be brought forth by the resurection of the dead through the triumph and glory of the lamb who was slain who was in the bosom of the father before the worlds were made
				% \item[BC] 
				% \item[]
			\end{compactdesc}
		\end{framed}

	% -----
	\subsubsection{Section 76:40} % *DC76-40 
	% -----
		\label{DC76-40}

		And this is the gospel,%
			\footnote{%
				% \label{}%
				We're about to get one of the scriptural definitions of the term ``gospel.'' See some others \hyperref[GOSPEL-definition]{here}.
				}
		the glad tidings, which the voice out of the heavens bore record unto us---

		% \begin{framed}
		% \scriptsize
			% \begin{compactdesc}
				% \item[JSP] 
				% % \item[BC] 
				% % \item[]
			% \end{compactdesc}
		% \end{framed}

	% -----
	\subsubsection{Section 76:41} % *DC76-41 
	% -----
		\label{DC76-41}

		That he came into the world, even Jesus, to be crucified for the world, and to bear the sins of the world, and to sanctify the world, and to cleanse it from all unrighteousness;

		% \begin{framed}
		% \scriptsize
			% \begin{compactdesc}
				% \item[JSP] 
				% % \item[BC] 
				% % \item[]
			% \end{compactdesc}
		% \end{framed}

	% -----
	\subsubsection{Section 76:42} % *DC76-42 
	% -----
		\label{DC76-42}

		That through him all might be saved whom the Father had put into his power and made by him;

		% \begin{framed}
		% \scriptsize
			% \begin{compactdesc}
				% \item[JSP] 
				% % \item[BC] 
				% % \item[]
			% \end{compactdesc}
		% \end{framed}

	% -----
	\subsubsection{Section 76:43} % *DC76-43 
	% -----
		\label{DC76-43}

		Who glorifies the Father, and saves all the works of his hands, except those sons of perdition who deny the Son after the Father has revealed him.%
			\footnote{%
				% \label{}%
				I haven't yet found any good scriptures that clearly discuss the idea of the Father revealing the Son to someone; see \hyperref[MATTHEW-11-27]{Matthew 11:27} and \hyperref[LUKE-10-22]{Luke 10:22} for Christ saying that he reveals the Father to people. On the other hand, there are a number of instances where the Father \textit{introduces} the Son, and in that sense reveals him to be the Son, and maybe that's what this is referring to here; that would be places like \hyperref[MATTHEW-3-17]{Matthew 3:17}, \hyperref[MATTHEW-17-5]{Matthew 17:5}, \hyperref[3NEPHI-11-7]{3 Nephi 11:7}, \href{https://www.josephsmithpapers.org/paper-summary/history-circa-june-1839-circa-1841-draft-2/3}{this account} of JS's first vision, and so on. So is it, then, that the sons of perdition have to have Christ revealed to them by the Father like it happens in these accounts? I'm not sure. See also \hyperref[2NEPHI-32-6]{2 Nephi 32:6}.
				}

		% \begin{framed}
		% \scriptsize
			% \begin{compactdesc}
				% \item[JSP] 
				% % \item[BC] 
				% % \item[]
			% \end{compactdesc}
		% \end{framed}

	% -----
	\subsubsection{Section 76:44} % *DC76-44 
	% -----
		\label{DC76-44}

		Wherefore, he saves all except them---they shall go away into everlasting punishment, which is endless punishment, which is eternal punishment,%
			\footnote{%
				% \label{}%
				See \hyperref[DC19-6]{DC 19:6ff}.
				}
		to reign with the devil and his angels in eternity, where their worm dieth not, and the fire is not quenched, which is their torment---

		\begin{framed}
		\scriptsize
			\begin{compactdesc}
				\item[JSP] wherefore he saveth all save them, and these shall go away into everlasting punishment which is eternal punishment to reign with the devel and his angels throughout all eternity where the worm dieth not and the fire is not quenched which is there torment
				% \item[BC] 
				% \item[]
			\end{compactdesc}
		\end{framed}

	% -----
	\subsubsection{Section 76:45} % *DC76-45 
	% -----
		\label{DC76-45}

		And the \hyperref[END]{end} thereof, neither the place thereof, nor their torment, no man knows;

		% \begin{framed}
		% \scriptsize
			% \begin{compactdesc}
				% \item[JSP] 
				% % \item[BC] 
				% % \item[]
			% \end{compactdesc}
		% \end{framed}

	% -----
	\subsubsection{Section 76:46} % *DC76-46 
	% -----
		\label{DC76-46}

		Neither was it revealed, neither is, neither will be revealed unto man, except to them who are made partakers thereof;

		% \begin{framed}
		% \scriptsize
			% \begin{compactdesc}
				% \item[JSP] 
				% % \item[BC] 
				% % \item[]
			% \end{compactdesc}
		% \end{framed}

	% -----
	\subsubsection{Section 76:47} % *DC76-47 
	% -----
		\label{DC76-47}

		Nevertheless, I, the Lord, show it by vision unto many, but straightway shut it up again;

		% \begin{framed}
		% \scriptsize
			% \begin{compactdesc}
				% \item[JSP] 
				% % \item[BC] 
				% % \item[]
			% \end{compactdesc}
		% \end{framed}

	% -----
	\subsubsection{Section 76:48} % *DC76-48 
	% -----
		\label{DC76-48}

		Wherefore, the \hyperref[END]{end}, the width, the height, the depth, and the misery thereof, they understand not, neither any man except those who are ordained unto this condemnation.%
			\footnote{%
				% \label{}%
				The ``ordained unto this condemnation'' language comes from \hyperref[JUDE-1-4]{Jude 1:4}. In Greek it is \gr{οἱ πάλαι προγεγραμμένοι εἰς τοῦτο τὸ κρίμα}; see \hyperref[KRIMA]{\grl{κρίμα}}. The short version of the Jude verse is that ``ordained'' may convey too strong a sense of an ordinance or judgment that was made regarding the condemnation---the word is actually \gr{προγράφω}, and BDAG 769a says that it's meaning \#1.b, ``what is written before, is found in an older document (by another author, as well).'' That entry glosses the Jude use as ``\textit{who for a long time have been marked out} (or \textit{written about}) \textit{for this judgment} (described in what follows).'' But note also that this could be a case of JS seeing the ``ordain'' language, and, not knowing the Greek, he uses his own sense of ``ordain'' to reveal/develop some theology around this point that goes beyond what the NT use would license. This is an important thing to keep in mind when studying JS's use of the NT.
				}

		% \begin{framed}
		% \scriptsize
			% \begin{compactdesc}
				% \item[JSP] 
				% % \item[BC] 
				% % \item[]
			% \end{compactdesc}
		% \end{framed}

	% -----
	\subsubsection{Section 76:49} % *DC76-49 
	% -----
		\label{DC76-49}

		And we heard the voice, saying: Write the vision, for lo, this is the end of the vision of the sufferings of the ungodly.%
			\footnote{%
				% \label{}%
				Strangely, ``ungodly'' shows up a bunch of times in Jude as well, including in verse \hyperref[JUDE-1-4]{Jude 1:4}, which is important to the previous verse in this section of the DC. ``Ungodly'' is also in \hyperref[JUDE-1-15]{verse 15} an amazing four times, and \hyperref[JUDE-1-18]{Jude 1:18} as well (see also \hyperref[2PETER-2-5]{2 Peter 2:5--6} and \hyperref[2PETER-3-7]{2 Peter 3:7}). In the Jude verses, ``ungodly'' is \hyperref[ASEBES]{\grl{ἀσεβής}}.
				}

		\begin{framed}
		\scriptsize
			\begin{compactdesc}
				\item[JSP] and we heard the voice saying write the vision for lo this is the end of the vision of the eternal sufferings of the undodly [ungodly]
				% \item[BC] 
				% \item[]
			\end{compactdesc}
		\end{framed}

	% -----
	\subsubsection{Section 76:50} % *DC76-50 
	% -----
		\label{DC76-50}

			\footnote{%
				% \label{}%
				Hard shift at this point in the section---even the earliest published versions (EMS, 1835 DC, 1844 DC), which don't have verses or anything, have a paragraph break here. We have just finished speaking about a group of people who will not be resurrected; the vision now moves on to describe people who \textit{are} resurrected, and in particular, those which are resurrected in the ``resurrection of the just,'' which is mentioned in \hyperref[DC76-17]{verse 17}.
				}%
		And again we bear record---for we saw and heard, and this is the testimony of the gospel of Christ concerning them who shall come forth in the resurrection of the just---%
			\footnote{%
				% \label{}%
				See notes at \hyperref[DC76-17]{verse 17} for the ``resurrection of the just.''
				}

		% \begin{framed}
		% \scriptsize
			% \begin{compactdesc}
				% \item[JSP] 
				% % \item[BC] 
				% % \item[]
			% \end{compactdesc}
		% \end{framed}

	% -----
	\subsubsection{Section 76:51} % *DC76-51 
	% -----
		\label{DC76-51}

		They are they who received the testimony of Jesus, and believed on his name and were baptized after the manner of his burial, being buried in the water in his name, and this according to the commandment which he has given---

		% \begin{framed}
		% \scriptsize
			% \begin{compactdesc}
				% \item[JSP] 
				% % \item[BC] 
				% % \item[]
			% \end{compactdesc}
		% \end{framed}

	% -----
	\subsubsection{Section 76:52} % *DC76-52 
	% -----
		\label{DC76-52}

		That by keeping the commandments they might be washed and cleansed from all their sins, and receive the Holy Spirit by the laying on of the hands of him who is ordained and sealed unto this power;

		% \begin{framed}
		% \scriptsize
			% \begin{compactdesc}
				% \item[JSP] 
				% % \item[BC] 
				% % \item[]
			% \end{compactdesc}
		% \end{framed}

	% -----
	\subsubsection{Section 76:53} % *DC76-53 
	% -----
		\label{DC76-53}

		And who overcome by faith,%
			\footnote{%
				% \label{}%
				See \hyperref[1JOHN-5-4]{1 John 5:4--5}: ``For whatsoever is born of God overcometh the world: and this is the victory that overcometh the world, even our faith. Who is he that overcometh the world, but he that believeth that Jesus is the Son of God?'' See notes there, and also see \hyperref[JOHN-16-33]{John 16:33}. Faith is the force which produces our victory over the world, as is made clear here in this DC verse.
				}
		and are sealed%
			\footnote{%
				% \label{}%
				The whole ``sealed by the Holy Spirit of promise'' language comes from \hyperref[EPHESIANS-1-13]{Ephesians 1:13}; see my notes there, and at \hyperref[HSOP]{Holy Spirit of Promise} (\hyperref[DC132]{DC 132} talks a lot about this as well). The Greek word for ``sealed'' is \hyperref[SPHRAGIZO]{\grl{σφραγίζω}}.
				}
		by the Holy Spirit of promise,%
			\footnote{%
				% \label{}%
				See \hyperref[HSOP]{Holy Spirit of Promise}, and especially the \hyperref[HSOP-DC]{DC uses} of the term.
				}
		which the Father sheds forth%
			\footnote{%
				% \label{}%
				An unusual construction, used in various places; see \hyperref[SHED-abroad]{this study} on ``shed abroad'' and ``shed forth.''
				}
		upon all those who are just and true.%
			\footnote{%
				% \label{}%
				The ``just and true'' combination appears in various interesting places in the scriptures---see \hyperref[REVELATIONNT-15-3]{Revelation 15:3} (\gr{δίκαιαι καὶ ἀληθιναὶ αἱ ὁδοί σου}), \hyperref[1NEPHI-14-23]{1 Nephi 14:23}, \hyperref[MOSIAH-4-12]{Mosiah 4:12}, \hyperref[MORONI-10-6]{Moroni 10:6}, and \hyperref[DC20-30]{DC 20:30--31}.
				}

		% \begin{framed}
		% \scriptsize
			% \begin{compactdesc}
				% \item[JSP] 
				% % \item[BC] 
				% % \item[]
			% \end{compactdesc}
		% \end{framed}

	% -----
	\subsubsection{Section 76:54} % *DC76-54 
	% -----
		\label{DC76-54}

		They are they who are the church of the Firstborn.%
			\footnote{%
				% \label{}%
				See \hyperref[COLOSSIANS-1-18]{Colossians 1:18} but especially \hyperref[HEBREWS-12-23]{Hebrews 12:23}, which is where this ``church of the Firstborn'' language comes from. See \hyperref[CHURCH-PHRASES]{this list} for other uses of the phrase ``church of the Firstborn,'' and see my notes on the term at the Hebrews verse. The key insight is that ``firstborn'' in the Greek is plural, not singular, and so that term in the phrase does not refer to Christ.
				}

		% \begin{framed}
		% \scriptsize
			% \begin{compactdesc}
				% \item[JSP] 
				% % \item[BC] 
				% % \item[]
			% \end{compactdesc}
		% \end{framed}

	% -----
	\subsubsection{Section 76:55} % *DC76-55 
	% -----
		\label{DC76-55}

		They are they into whose hands the Father has given all things---%
			\footnote{%
				% \label{}%
				I think the best explanation of the idea in this verse is in \hyperref[DC50-26]{DC 50:26--28}; 27 reads, ``Wherefore, he is possessor of all things; for all things are subject unto him, both in heaven and on the earth, the life and the light, the Spirit and the power, sent forth by the will of the Father through Jesus Christ, his Son.'' See the notes at that verse. See also \hyperref[JOHN-3-35]{John 3:35}, \hyperref[JOHN-13-3]{John 13:3}, and \hyperref[DC76-59]{verse 59} in this section. In thinking about the ``Father'' specifically doing this, it reminds me of \hyperref[2NEPHI-31-20]{2 Nephi 31:20}: ``behold, thus saith the Father, ye shall have eternal life.'' 
				}

		% \begin{framed}
		% \scriptsize
			% \begin{compactdesc}
				% \item[JSP] 
				% % \item[BC] 
				% % \item[]
			% \end{compactdesc}
		% \end{framed}

	% -----
	\subsubsection{Section 76:56} % *DC76-56 
	% -----
		\label{DC76-56}

		They are they who are priests and kings,%
			\footnote{%
				% \label{}%
				See \hyperref[GENESIS-14-18]{Genesis 14:18}, \hyperref[EXODUS-19-6]{Exodus 19:6}, \hyperref[HEBREWS-7-1]{Hebrews 7:1}, \hyperref[1PETER-2-9]{1 Peter 2:9}, and \hyperref[REVELATIONNT-1-6]{Revelation 1:6} and \hyperref[REVELATIONNT-5-10]{5:10}. This language is in the \hyperref[TEMPLE-ENDOWMENT-ENDOW-PRE-1990-INTRO]{introduction} to the endowment: ``Brethren, you have been washed and pronounced clean, or that through your faithfulness you may become clean, from the blood and sins of this generation. You have been anointed to become hereafter kings and priests unto the most high God, to rule and reign in the house of Israel forever. Sisters, you have been washed and anointed to become queens and priestesses to your husbands. Brethren and sisters, if you are true and faithful, the day will come when you will be chosen, called up, and anointed kings and queens, priests and priestesses, whereas you are now anointed only to become such. The realization of these blessings depends upon your faithfulness.''
				}
		who have received of his fulness,%
			\footnote{%
				% \label{}%
				See \hyperref[DC76-20]{verse 20}, \hyperref[DC76-71]{verse 71}, and \hyperref[DC76-77]{verse 77}; see also \hyperref[DC93-4]{DC 93:4} and \hyperref[DC93-16]{16--20}. The idea of receiving the ``fulness'' of the Father seems related to the previous verse, in 55, where the Father ``gives all things'' into the hands of this group of people.
				}
		and of his glory;

		% \begin{framed}
		% \scriptsize
			% \begin{compactdesc}
				% \item[JSP] 
				% % \item[BC] 
				% % \item[]
			% \end{compactdesc}
		% \end{framed}

	% -----
	\subsubsection{Section 76:57} % *DC76-57 
	% -----
		\label{DC76-57}

		And are priests of the Most High, after the order of Melchizedek, which was after the order of Enoch,%
			\footnote{%
				% \label{}%
				For the ``order of Enoch,'' see \hyperref[PRIESTHOOD-enoch]{here}.
				}
		which was after the order of the Only Begotten Son.%
			\footnote{%
				% \label{}%
				Fundamentally, the priesthood is ``\textit{the Holy Priesthood, after the Order of the Son of God}'' (\hyperref[DC107-3]{DC 107:3}; see also 1--6 and notes there). Sometimes it receives another name, depending on who is alive and the ``highest'' priest at the time (or the one who holds all the keys; see the note above on the ``order of Enoch''). But the priesthood is after the order of the Son of God, or Christ. This is why the fundamental law of the celestial kingdom is called the ``law of Christ'' (more on that phrase \hyperref[LAW-PHRASES]{here}; see also \hyperref[DC88-21]{DC 88:21}).
				}

		% \begin{framed}
		% \scriptsize
			% \begin{compactdesc}
				% \item[JSP] 
				% % \item[BC] 
				% % \item[]
			% \end{compactdesc}
		% \end{framed}

	% -----
	\subsubsection{Section 76:58} % *DC76-58 
	% -----
		\label{DC76-58}

		Wherefore, as it is written, they are gods, even the sons of God---%
			\footnote{%
				% \label{}%
				Apparently a reference to \hyperref[PSALMS-82-6]{Psalms 82:6}.
				}

		% \begin{framed}
		% \scriptsize
			% \begin{compactdesc}
				% \item[JSP] 
				% % \item[BC] 
				% % \item[]
			% \end{compactdesc}
		% \end{framed}

	% -----
	\subsubsection{Section 76:59} % *DC76-59 
	% -----
		\label{DC76-59}

		Wherefore, all things are theirs, whether life or death, or things present, or things to come, all are theirs and they are Christ's,%
			\footnote{%
				% \label{}%
				See \hyperref[MARK-9-41]{Mark 9:41}; \hyperref[DC101-3]{DC 101:3}; 
				}
		and Christ is God's.

		% \begin{framed}
		% \scriptsize
			% \begin{compactdesc}
				% \item[JSP] 
				% % \item[BC] 
				% % \item[]
			% \end{compactdesc}
		% \end{framed}

	% -----
	\subsubsection{Section 76:60} % *DC76-60 
	% -----
		\label{DC76-60}

		And they shall overcome%
			\footnote{%
				% \label{}%
				The Greek would be \hyperref[NIKAO]{\grl{νικάω}}.
				}
		all things.

		% \begin{framed}
		% \scriptsize
			% \begin{compactdesc}
				% \item[JSP] 
				% % \item[BC] 
				% % \item[]
			% \end{compactdesc}
		% \end{framed}

	% -----
	\subsubsection{Section 76:61} % *DC76-61 
	% -----
		\label{DC76-61}

		Wherefore, let no man glory in man, but rather let him glory in God, who shall subdue all enemies%
			\footnote{%
				% \label{}%
				For the idea of overcoming ``enemies,'' see \hyperref[SALVATION-JS]{JS on ``salvation''}.
				}
		under his feet.%
			\footnote{%
				% \label{}%
				It's surprising to me that the ``subdue enemies under feet'' construction is unique to restoration scripture; a search of ``subdu* enem* feet*'' only gives you this verse, \hyperref[DC76-106]{verse 106}, and \hyperref[DC58-22]{DC 58:22}. If you search ``subdu* enem*'' you get \hyperref[1CHRONICLES-17-10]{1 Chronicles 17:10} and \hyperref[PSALMS-81-14]{Psalms 81:14}, plus \hyperref[DC65-6]{DC 65:6}, but that's it. 
				}

		% \begin{framed}
		% \scriptsize
			% \begin{compactdesc}
				% \item[JSP] 
				% % \item[BC] 
				% % \item[]
			% \end{compactdesc}
		% \end{framed}

	% -----
	\subsubsection{Section 76:62} % *DC76-62 
	% -----
		\label{DC76-62}

		These shall dwell in the presence of God and his Christ forever and ever.

		% \begin{framed}
		% \scriptsize
			% \begin{compactdesc}
				% \item[JSP] 
				% % \item[BC] 
				% % \item[]
			% \end{compactdesc}
		% \end{framed}

	% -----
	\subsubsection{Section 76:63} % *DC76-63 
	% -----
		\label{DC76-63}

		These are they whom he shall bring with him, when he shall come in the clouds of heaven to reign on the earth over his people.

		% \begin{framed}
		% \scriptsize
			% \begin{compactdesc}
				% \item[JSP] 
				% % \item[BC] 
				% % \item[]
			% \end{compactdesc}
		% \end{framed}

	% -----
	\subsubsection{Section 76:64} % *DC76-64 
	% -----
		\label{DC76-64}

		These are they who shall have part in the first resurrection.

		% \begin{framed}
		% \scriptsize
			% \begin{compactdesc}
				% \item[JSP] 
				% % \item[BC] 
				% % \item[]
			% \end{compactdesc}
		% \end{framed}

	% -----
	\subsubsection{Section 76:65} % *DC76-65 
	% -----
		\label{DC76-65}

		These are they who shall come forth in the resurrection of the just.%
			\footnote{%
				% \label{}%
				The way this verse and the previous are set up makes it seem as though they are drawing a distinction between the ``first resurrection'' and the ``resurrection of the just.'' My impression from other verses is that the latter would be narrower and take place during the period of the former, but I haven't studied this enough yet.
				}

		% \begin{framed}
		% \scriptsize
			% \begin{compactdesc}
				% \item[JSP] 
				% % \item[BC] 
				% % \item[]
			% \end{compactdesc}
		% \end{framed}

	% -----
	\subsubsection{Section 76:66} % *DC76-66 
	% -----
		\label{DC76-66}

			\footnote{%
				% \label{}%
				This verse and the next three return again to \hyperref[HEBREWS-12-22]{Hebrews 12:22--24}. It does add a bit to those verses, however---``the heavenly place, the holiest of all'' and ``church of Enoch'' aren't in the Hebrews verses.
				}%
		These are they who are come unto Mount Zion, and unto the city of the living God, the heavenly place, the holiest of all.

		% \begin{framed}
		% \scriptsize
			% \begin{compactdesc}
				% \item[JSP] 
				% % \item[BC] 
				% % \item[]
			% \end{compactdesc}
		% \end{framed}

	% -----
	\subsubsection{Section 76:67} % *DC76-67 
	% -----
		\label{DC76-67}

		These are they who have come to an innumerable company of angels, to the general assembly and church of Enoch,%
			\footnote{%
				% \label{}%
				This is the only occurrence of ``church of Enoch'' in all the scriptures.
				}
		and of the Firstborn.

		% \begin{framed}
		% \scriptsize
			% \begin{compactdesc}
				% \item[JSP] 
				% % \item[BC] 
				% % \item[]
			% \end{compactdesc}
		% \end{framed}

	% -----
	\subsubsection{Section 76:68} % *DC76-68 
	% -----
		\label{DC76-68}

		These are they whose names are written in heaven, where God and Christ are the judge of all.

		% \begin{framed}
		% \scriptsize
			% \begin{compactdesc}
				% \item[JSP] 
				% % \item[BC] 
				% % \item[]
			% \end{compactdesc}
		% \end{framed}

	% -----
	\subsubsection{Section 76:69} % *DC76-69 
	% -----
		\label{DC76-69}

		These are they who are just men made perfect through Jesus the mediator of the new covenant,%
			\footnote{%
				% \label{}%
				This verse was my favorite scripture when I was a missionary, and in particular this first half. It's inspiring to know that we just need to be \textit{just}, or \hyperref[DIKAIOS]{\grl{δίκαιος}}, and that's enough. We are then made perfect through Jesus the mediator. Some of the language comes from \hyperref[HEBREWS-12-23]{Hebrews 12:23--24}; the ``made perfect'' there is \gr{τετελειωμένων}, a participle from \hyperref[TELEIOO]{\grl{τελειόω}}. Note that one of the definition \#2.b of \gr{τελειόω} is ``bring to full measure, fill the measure of,'' which is also language used in restoration scripture, such as \hyperref[DC88-19]{DC 88:19} and the endowment ceremony.
				}
		who wrought out this perfect atonement%
			\footnote{%
				% \label{}%
				The only use of ``perfect atonement'' in the scriptures. It's interesting that it would be called ``perfect'' here, since we are talking about just people being made \textit{perfect} through the \textit{perfect} atonement. It reminds me of \hyperref[DC19-19]{DC 19:19}, where Christ ``finished [his] preparations unto the children of men.'' Meaning \#1 of \hyperref[TELEIOO]{\grl{τελειόω}}---see the note for the first part of this verse---is ``complete, bring to an end, finish, accomplish.'' Christ \gr{τελειόω}'d the atonement, and in doing so, \gr{τελειόω}'d us as well.
				}
		through the shedding of his own blood.

		% \begin{framed}
		% \scriptsize
			% \begin{compactdesc}
				% \item[JSP] 
				% % \item[BC] 
				% % \item[]
			% \end{compactdesc}
		% \end{framed}

	% -----
	\subsubsection{Section 76:70} % *DC76-70 
	% -----
		\label{DC76-70}

			\footnote{%
				% \label{}%
				This verse begins a new sub-section which drops the language from Hebrews and instead picks up \hyperref[1CORINTHIANS-15-40]{1 Corinthians 15:40}.
				}%
		These are they whose bodies are celestial, whose%
			\footnote{%
				% \label{}%
				The antecedent of this ``whose'' seems to be the ``they'' earlier in the verse, rather than ``bodies.'' 
				}
		glory is that of the sun, even the glory of God, the highest of all, whose glory the sun%
		\footnote{%
			% \label{}%
			Always makes me think of the \hyperref[TTW-DEVELOPMENT-PHILOSOPHICAL-001]{allegory of the cave} from \textit{Republic} book VII.
			} 
		of the firmament is written of as being typical.

		% \begin{framed}
		% \scriptsize
			% \begin{compactdesc}
				% \item[JSP] 
				% % \item[BC] 
				% % \item[]
			% \end{compactdesc}
		% \end{framed}

	% -----
	\subsubsection{Section 76:71} % *DC76-71 
	% -----
		\label{DC76-71}

			\footnote{%
				% \label{}%
				Now a much larger break---transitioning from the celestial world to the terrestrial world.
				}%
		And again, we saw the terrestrial world, and behold and lo, these are they who are of the terrestrial, whose glory differs from that of the church of the Firstborn who have received the fulness of the Father, even as that of the moon differs from the sun in the firmament.

		% \begin{framed}
		% \scriptsize
			% \begin{compactdesc}
				% \item[JSP] 
				% % \item[BC] 
				% % \item[]
			% \end{compactdesc}
		% \end{framed}

	% -----
	\subsubsection{Section 76:72} % *DC76-72 
	% -----
		\label{DC76-72}

		Behold, these are they who died without law;%
			\footnote{%
				% \label{}%
				Possibly a reference to \hyperref[ROMANS-7-8]{Romans 7:8--9}. See the \hyperref[ATONEMENT-categories]{categories} of people to whom the atonement applies. Another way to view this comment is to see it as a reference to the priesthood, or the \textit{lack} of integration of someone's life into the laws of the priesthood that govern salvation. See also \hyperref[ALMA-42-17]{Alma 42:17--22}.
				}

		% \begin{framed}
		% \scriptsize
			% \begin{compactdesc}
				% \item[JSP] 
				% % \item[BC] 
				% % \item[]
			% \end{compactdesc}
		% \end{framed}

	% -----
	\subsubsection{Section 76:73} % *DC76-73 
	% -----
		\label{DC76-73}

		And also they who are the spirits of men kept in prison, whom the Son visited, and preached the gospel unto them,%
			\footnote{%
				% \label{}%
				See \hyperref[1PETER-3-18]{1 Peter 3:18--20}, \hyperref[1PETER-4-6]{1 Peter 4:6}, and \hyperref[DC138]{DC 138}.
				}
		that they might be judged according to men in the flesh;%
			\footnote{%
				% \label{}%
				This entire last part, ``that they might be judged according to men in the flesh,'' comes from \hyperref[1PETER-4-6]{1 Peter 4:6}; see \hyperref[FLESH-PHRASES]{here} for more occurrences of ``according to men in the flesh.''
				}

		% \begin{framed}
		% \scriptsize
			% \begin{compactdesc}
				% \item[JSP] 
				% % \item[BC] 
				% % \item[]
			% \end{compactdesc}
		% \end{framed}

	% -----
	\subsubsection{Section 76:74} % *DC76-74 
	% -----
		\label{DC76-74}

		Who received not the testimony of Jesus in the flesh, but afterwards received it.

		% \begin{framed}
		% \scriptsize
			% \begin{compactdesc}
				% \item[JSP] 
				% % \item[BC] 
				% % \item[]
			% \end{compactdesc}
		% \end{framed}

	% -----
	\subsubsection{Section 76:75} % *DC76-75 
	% -----
		\label{DC76-75}

		These are they who are honorable men of the earth, who were blinded by the craftiness of men.%
			\footnote{%
				% \label{}%
				Possibly a reference to \hyperref[EPHESIANS-4-14]{Ephesians 4:14}. In several places the BOM talks about ``cunning and craftiness,'' and the Ephesians verse uses ``cunning craftiness.'' 
				}

		% \begin{framed}
		% \scriptsize
			% \begin{compactdesc}
				% \item[JSP] 
				% % \item[BC] 
				% % \item[]
			% \end{compactdesc}
		% \end{framed}

	% -----
	\subsubsection{Section 76:76} % *DC76-76 
	% -----
		\label{DC76-76}

		These are they who receive of his glory, but not of his fulness.%
			\footnote{%
				% \label{}%
				Compare \hyperref[DC88-28]{DC 88:28--32}.
				}

		% \begin{framed}
		% \scriptsize
			% \begin{compactdesc}
				% \item[JSP] 
				% % \item[BC] 
				% % \item[]
			% \end{compactdesc}
		% \end{framed}

	% -----
	\subsubsection{Section 76:77} % *DC76-77 
	% -----
		\label{DC76-77}

		These are they who receive of the presence of the Son, but not of the fulness of the Father.

		% \begin{framed}
		% \scriptsize
			% \begin{compactdesc}
				% \item[JSP] 
				% % \item[BC] 
				% % \item[]
			% \end{compactdesc}
		% \end{framed}

	% -----
	\subsubsection{Section 76:78} % *DC76-78 
	% -----
		\label{DC76-78}

		Wherefore, they are bodies terrestrial, and not bodies celestial, and differ in glory as the moon differs from the sun.

		% \begin{framed}
		% \scriptsize
			% \begin{compactdesc}
				% \item[JSP] 
				% % \item[BC] 
				% % \item[]
			% \end{compactdesc}
		% \end{framed}

	% -----
	\subsubsection{Section 76:79} % *DC76-79 
	% -----
		\label{DC76-79}

		These are they who are not \hyperref[VALIANT]{valiant} in the testimony of Jesus; wherefore, they obtain not the crown over the kingdom of our God.

		% \begin{framed}
		% \scriptsize
			% \begin{compactdesc}
				% \item[JSP] 
				% % \item[BC] 
				% % \item[]
			% \end{compactdesc}
		% \end{framed}

	% -----
	\subsubsection{Section 76:80} % *DC76-80 
	% -----
		\label{DC76-80}

		And now this is the end of the vision which we saw of the terrestrial, that the Lord commanded us to write while we were yet in the Spirit.

		% \begin{framed}
		% \scriptsize
			% \begin{compactdesc}
				% \item[JSP] 
				% % \item[BC] 
				% % \item[]
			% \end{compactdesc}
		% \end{framed}

	% -----
	\subsubsection{Section 76:81} % *DC76-81 
	% -----
		\label{DC76-81}

		And again, we saw the glory of the \hyperref[TELESTIAL]{telestial},%
			\footnote{%
				% \label{}%
				This is the first use of the term ``telestial'' in restoration scripture, and it occurs only in DC 76 and in \hyperref[DC88]{DC 88}, nowhere else. It also appears in a few places in the \hyperref[TELESTIAL-TEMPLE]{endowment ceremony}.
				}
		which glory is that of the lesser, even as the glory of the stars differs from that of the glory of the moon in the firmament.

		% \begin{framed}
		% \scriptsize
			% \begin{compactdesc}
				% \item[JSP] 
				% % \item[BC] 
				% % \item[]
			% \end{compactdesc}
		% \end{framed}

	% -----
	\subsubsection{Section 76:82} % *DC76-82 
	% -----
		\label{DC76-82}

		These are they who received not the gospel of Christ, neither the testimony of Jesus.%
			\footnote{%
				% \label{}%
				Notice that the ``gospel of Christ'' and ``the testimony of Jesus'' are spoken of as separate things in this verse---you can receive one without receiving the other. This very section speaks in particular about the gospel in various places, and connects it somewhat directly with going beyond the testimony of Jesus into receiving ordinances; see \hyperref[DC76-51]{verse 51} for example, and \hyperref[DC76-101]{verse 101}. A related term is the ``doctrine of Christ'' as spoken of in \hyperref[2NEPHI-31-21]{2 Nephi 31:21}, which is a summation of everything that Nephi taught about baptism and ordinances in that chapter. Sherem connects the terms at \hyperref[JACOB-7-6]{Jacob 7:6}.
				}

		% \begin{framed}
		% \scriptsize
			% \begin{compactdesc}
				% \item[JSP] 
				% % \item[BC] 
				% % \item[]
			% \end{compactdesc}
		% \end{framed}

	% -----
	\subsubsection{Section 76:83} % *DC76-83 
	% -----
		\label{DC76-83}

		These are they who deny not the Holy Spirit.

		% \begin{framed}
		% \scriptsize
			% \begin{compactdesc}
				% \item[JSP] 
				% % \item[BC] 
				% % \item[]
			% \end{compactdesc}
		% \end{framed}

	% -----
	\subsubsection{Section 76:84} % *DC76-84 
	% -----
		\label{DC76-84}

		These are they who are thrust down to hell.

		% \begin{framed}
		% \scriptsize
			% \begin{compactdesc}
				% \item[JSP] 
				% % \item[BC] 
				% % \item[]
			% \end{compactdesc}
		% \end{framed}

	% -----
	\subsubsection{Section 76:85} % *DC76-85 
	% -----
		\label{DC76-85}

		These are they who shall not be redeemed from the devil until the last resurrection, until the Lord, even Christ the Lamb, shall have finished his work.

		% \begin{framed}
		% \scriptsize
			% \begin{compactdesc}
				% \item[JSP] 
				% % \item[BC] 
				% % \item[]
			% \end{compactdesc}
		% \end{framed}

	% -----
	\subsubsection{Section 76:86} % *DC76-86 
	% -----
		\label{DC76-86}

		These are they who receive not of his fulness in the eternal world, but of the Holy Spirit through the ministration of the terrestrial;

		% \begin{framed}
		% \scriptsize
			% \begin{compactdesc}
				% \item[JSP] 
				% % \item[BC] 
				% % \item[]
			% \end{compactdesc}
		% \end{framed}

	% -----
	\subsubsection{Section 76:87} % *DC76-87 
	% -----
		\label{DC76-87}

		And the terrestrial through the ministration of the celestial.

		% \begin{framed}
		% \scriptsize
			% \begin{compactdesc}
				% \item[JSP] 
				% % \item[BC] 
				% % \item[]
			% \end{compactdesc}
		% \end{framed}

	% -----
	\subsubsection{Section 76:88} % *DC76-88 
	% -----
		\label{DC76-88}

		And also the telestial receive it of the administering of angels who are appointed to minister for them, or who are appointed to be ministering spirits for them;%
			\footnote{%
				% \label{}%
				\hyperref[DC107-20]{DC 107:20} tells us that the Aaronic priesthood holds ``the keys of the ministering of angels.'' In the endowment ceremony, you receive the first token of the Aaronic priesthood while still in the garden; immediately after you enter into the telestial world. In that world you then receive the second token of the Aaronic priesthood, and you put on the robes in the Aaronic way before changing them to the Melchizedek way, still in the telestial world. Notice that Peter, James, and John visit you, as ministers or as ministering angels, while in the telestial room, and while you are receiving these tokens of the Aaronic priesthood, in line with what is described here in this verse. Immediately after changing the robes to the right shoulder, you enter the terrestrial room. In that room you receive both tokens of the Melchizedek priesthood---note that both are clear references to Christ, along the lines of \hyperref[DC76-77]{verse 77} above, being in the terrestrial world. Once you'e received all the tokens of the Melchizedek priesthood, and you get the name of the second token at the veil, you are ready to enter into the presence of Elohim, in the celestial kingdom: ``Brethren and sisters, we are instructed to introduce you at the veil, where you will receive the name of the second token of the Melchizedek priesthood, the patriarchal grip, or Sure Sign of the Nail, preparatory to your entering into the presence of the Lord.'' This is a very beautiful connection between tihs revelation, which comes ten years before the restoration of the endowment, and the endowment ceremony itself. 
				}
		for they shall be heirs of salvation.

		% \begin{framed}
		% \scriptsize
			% \begin{compactdesc}
				% \item[JSP] 
				% % \item[BC] 
				% % \item[]
			% \end{compactdesc}
		% \end{framed}

	% -----
	\subsubsection{Section 76:89} % *DC76-89 
	% -----
		\label{DC76-89}

		And thus we saw, in the heavenly vision, the glory of the telestial, which surpasses%
		\footnote{%
			% \label{}%
			The word ``surpass'' only occurs twice in all the scriptures---here, and in \hyperref[DC76-114]{verse 114}.
			} 
		all understanding;%
			\footnote{%
				% \label{}%
				See BY, \hyperref[EMS-1865-BY-8-OCT-1865]{8 October 1865}: ``what will become of the inhabitants of the earth that have and now live all they that don’t sin against the Holy Ghost will be saved in a kingdom somewhere and they fare better than ever imagined and Father Wesley get a greater glory [than] ever entered into his heart but [when/one/won't?] go where the Father and Son are and so it will be with others.''
				}

		% \begin{framed}
		% \scriptsize
			% \begin{compactdesc}
				% \item[JSP] 
				% % \item[BC] 
				% % \item[]
			% \end{compactdesc}
		% \end{framed}

	% -----
	\subsubsection{Section 76:90} % *DC76-90 
	% -----
		\label{DC76-90}

		And no man knows it except him to whom God has revealed it.

		% \begin{framed}
		% \scriptsize
			% \begin{compactdesc}
				% \item[JSP] 
				% % \item[BC] 
				% % \item[]
			% \end{compactdesc}
		% \end{framed}

	% -----
	\subsubsection{Section 76:91} % *DC76-91 
	% -----
		\label{DC76-91}

		And thus we saw the glory of the terrestrial which excels in all things the glory of the telestial, even in glory, and in power, and in might, and in dominion.

		% \begin{framed}
		% \scriptsize
			% \begin{compactdesc}
				% \item[JSP] 
				% % \item[BC] 
				% % \item[]
			% \end{compactdesc}
		% \end{framed}

	% -----
	\subsubsection{Section 76:92} % *DC76-92 
	% -----
		\label{DC76-92}

		And thus we saw the glory of the celestial, which excels in all things---where God, even the Father, reigns upon his throne forever and ever;

		% \begin{framed}
		% \scriptsize
			% \begin{compactdesc}
				% \item[JSP] 
				% % \item[BC] 
				% % \item[]
			% \end{compactdesc}
		% \end{framed}

	% -----
	\subsubsection{Section 76:93} % *DC76-93 
	% -----
		\label{DC76-93}

		Before whose throne all things bow in humble reverence, and give him glory forever and ever.

		% \begin{framed}
		% \scriptsize
			% \begin{compactdesc}
				% \item[JSP] 
				% % \item[BC] 
				% % \item[]
			% \end{compactdesc}
		% \end{framed}

	% -----
	\subsubsection{Section 76:94} % *DC76-94 
	% -----
		\label{DC76-94}

		They who dwell in his presence are the church of the Firstborn; and they see as they are seen, and know as they are known,%
			\footnote{%
				% \label{}%
				I can't find a scriptural connection to the ``see as they are seen/known as they are known'' bit. Maybe I'm missing it somewhere but I can't find anything. The key question seems to be, seen and known by whom? The obvious answer is God, but given that everyone else receives the powers of God, it could also be that the verse is referring to each person who is part of the church of the Firstborn. Some scriptures to keep in mind: \hyperref[GALATIANS-4-9]{Galatians 4:9}, \hyperref[PHILIPPIANS-3-12]{Philippians 3:12} (uses ``apprehend/apprehended of Christ''), 
				}
		having received of his fulness and of his grace;

		% \begin{framed}
		% \scriptsize
			% \begin{compactdesc}
				% \item[JSP] 
				% % \item[BC] 
				% % \item[]
			% \end{compactdesc}
		% \end{framed}

	% -----
	\subsubsection{Section 76:95} % *DC76-95 
	% -----
		\label{DC76-95}

		And he makes them equal%
			\footnote{%
				% \label{}%
				There are at least two ways of taking this ``equal'' part---it could be that God makes these individuals equal \textit{to each other}, but it could also be that he makes them equal to \textit{himself} in these respects. \hyperref[MATTHEW-20-12]{Matthew 20:12} might suggest the first interpretation, and is an interesting connection to this verse; see also \hyperref[MATTHEW-10-25]{Matthew 10:25}, \hyperref[JOHN-5-18]{John 5:18} and \hyperref[DC88-107]{DC 88:107}.
				}
		in power, and in might, and in dominion.%
			\footnote{%
				% \label{}%
				This triple of ``power, might, dominion'' occurs in \hyperref[EPHESIANS-1-21]{Ephesians 1:21}; the ``power and might'' there are \gr{ἐξουσία} and \gr{δύναμις}. The combination of the three terms is also at \hyperref[1NEPHI-22-24]{1 Nephi 22:24}, \hyperref[ALMA-5-50]{Alma 5:50}, \hyperref[ALMA-12-15]{Alma 12:15}, and before in this section in \hyperref[DC76-91]{verse 91}; see also \hyperref[DC109-77]{DC 109:77}.
				}

		% \begin{framed}
		% \scriptsize
			% \begin{compactdesc}
				% \item[JSP] 
				% % \item[BC] 
				% % \item[]
			% \end{compactdesc}
		% \end{framed}

	% -----
	\subsubsection{Section 76:96} % *DC76-96 
	% -----
		\label{DC76-96}

		And the glory of the celestial is one,%
			\footnote{%
				% \label{}%
				This bit and the next two verses have material from \hyperref[1CORINTHIANS-15-40]{1 Corinthians 15:40--41}. This verse here, when you read it with the next two, make it seem as though the ``one'' means something like ``unitary''---``the glory of the celestian is unitary, one, single,'' like everyone has the same glory. I don't think the Greek can be read this way---the AB translation of the second part of 15:40 brings this out: ``the splendor of the heavenly bodies is one thing; that of the earthly is another.'' In other words, the Greek is trying to underscore the distinction between the glories or levels of ``splendor.'' These verses here in DC can be read that way, and probably should be read that way, but maybe they are suggesting other things too.
				}
		even as the glory of the sun is one.

		% \begin{framed}
		% \scriptsize
			% \begin{compactdesc}
				% \item[JSP] 
				% % \item[BC] 
				% % \item[]
			% \end{compactdesc}
		% \end{framed}

	% -----
	\subsubsection{Section 76:97} % *DC76-97 
	% -----
		\label{DC76-97}

		And the glory of the terrestrial is one, even as the glory of the moon is one.

		% \begin{framed}
		% \scriptsize
			% \begin{compactdesc}
				% \item[JSP] 
				% % \item[BC] 
				% % \item[]
			% \end{compactdesc}
		% \end{framed}

	% -----
	\subsubsection{Section 76:98} % *DC76-98 
	% -----
		\label{DC76-98}

		And the glory of the telestial is one, even as the glory of the stars is one; for as one star differs from another star in glory, even so differs one from another in glory in the telestial world;

		% \begin{framed}
		% \scriptsize
			% \begin{compactdesc}
				% \item[JSP] 
				% % \item[BC] 
				% % \item[]
			% \end{compactdesc}
		% \end{framed}

	% -----
	\subsubsection{Section 76:99} % *DC76-99 
	% -----
		\label{DC76-99}

		For these are they who are of Paul, and of Apollos, and of Cephas.%
			\footnote{%
				% \label{}%
				See \hyperref[1CORINTHIANS-1-12]{1 Corinthians 1:12} and \hyperref[1CORINTHIANS-3-22]{3:22}.
				}

		% \begin{framed}
		% \scriptsize
			% \begin{compactdesc}
				% \item[JSP] 
				% % \item[BC] 
				% % \item[]
			% \end{compactdesc}
		% \end{framed}

	% -----
	\subsubsection{Section 76:100} % *DC76-100 
	% -----
		\label{DC76-100}

		These are they who say they are some of one and some of another---some of Christ%
			\footnote{%
				% \label{}%
				Notice that Christ is one of the possibilities here---so even though some groups of people claim to be of Christ, they have not taken the actions that were required of them, as stated in the next verse.
				}
		and some of John, and some of Moses, and some of Elias, and some of Esaias, and some of Isaiah, and some of Enoch;

		% \begin{framed}
		% \scriptsize
			% \begin{compactdesc}
				% \item[JSP] 
				% % \item[BC] 
				% % \item[]
			% \end{compactdesc}
		% \end{framed}

	% -----
	\subsubsection{Section 76:101} % *DC76-101 
	% -----
		\label{DC76-101}

		But received not the gospel, neither the testimony of Jesus, neither the prophets, neither the everlasting covenant.

		% \begin{framed}
		% \scriptsize
			% \begin{compactdesc}
				% \item[JSP] 
				% % \item[BC] 
				% % \item[]
			% \end{compactdesc}
		% \end{framed}

	% -----
	\subsubsection{Section 76:102} % *DC76-102 
	% -----
		\label{DC76-102}

		Last of all, these all are they who will not be gathered with the saints, to be caught up unto the church of the Firstborn, and received into the cloud.

		% \begin{framed}
		% \scriptsize
			% \begin{compactdesc}
				% \item[JSP] 
				% % \item[BC] 
				% % \item[]
			% \end{compactdesc}
		% \end{framed}

	% -----
	\subsubsection{Section 76:103} % *DC76-103 
	% -----
		\label{DC76-103}

		These are they who are liars, and sorcerers, and adulterers, and whoremongers, and whosoever loves and makes a lie.%
			\footnote{%
				% \label{}%
				Compare \hyperref[REVELATIONNT-21-8]{Revelation 21:8}, \hyperref[REVELATIONNT-22-15]{22:15}, and \hyperref[DC63-17]{DC 63:17}.
				}

		% \begin{framed}
		% \scriptsize
			% \begin{compactdesc}
				% \item[JSP] 
				% % \item[BC] 
				% % \item[]
			% \end{compactdesc}
		% \end{framed}

	% -----
	\subsubsection{Section 76:104} % *DC76-104 
	% -----
		\label{DC76-104}

		These are they who suffer the wrath of God on earth.

		% \begin{framed}
		% \scriptsize
			% \begin{compactdesc}
				% \item[JSP] 
				% % \item[BC] 
				% % \item[]
			% \end{compactdesc}
		% \end{framed}

	% -----
	\subsubsection{Section 76:105} % *DC76-105 
	% -----
		\label{DC76-105}

		These are they who suffer the vengeance of eternal fire.%
			\footnote{%
				% \label{}%
				As always with this stuff about punishment, you have to begin with \hyperref[DC19]{DC 19}.
				}

		% \begin{framed}
		% \scriptsize
			% \begin{compactdesc}
				% \item[JSP] 
				% % \item[BC] 
				% % \item[]
			% \end{compactdesc}
		% \end{framed}

	% -----
	\subsubsection{Section 76:106} % *DC76-106 
	% -----
		\label{DC76-106}

		These are they who are cast down to hell and suffer the wrath of Almighty God, until the fulness of times, when Christ shall have subdued all enemies under his feet, and shall have perfected his work;%
			\footnote{%
				% \label{}%
				The gospel of John has a number of places where it talks about Christ and ``perfecting'' this work, but they use the word ``finish'' in the KJV---see \hyperref[JOHN-4-34]{John 4:34}, \hyperref[JOHN-5-36]{5:36}, and \hyperref[JOHN-17-4]{17:4}. All three use forms of \hyperref[TELEIOO]{\grl{τελειόω}}, so that's the way I want to interpret what is being said here in the DC.
				}

		% \begin{framed}
		% \scriptsize
			% \begin{compactdesc}
				% \item[JSP] 
				% % \item[BC] 
				% % \item[]
			% \end{compactdesc}
		% \end{framed}

	% -----
	\subsubsection{Section 76:107} % *DC76-107 
	% -----
		\label{DC76-107}

		When he shall deliver up the kingdom,%
			\footnote{%
				% \label{}%
				That part of existence governed by the law of Christ.
				}
		and present it unto the Father, spotless,%
			\footnote{%
				% \label{}%
				``Spotless'' does not occur in the Bible at all; it's a word introduced into the gospel lexicon by the BOM. I almost want to take the term here as meaning ``perfectly organized''---everything is in its place, obeying the laws it chooses to obey as well as it can, filling the measure of its creation. See my note at \hyperref[DC76-69]{verse 69}. The purpose of an implementation of the gospel on a world is to impose order and organization on all matter and life; see \hyperref[ENDOWMENT-MEASURE]{this comment} in the endowment ceremony. Maybe another way of putting it is to say, ``spotless'' only has a meaning relative to a set of laws; it doesn't really mean anything by itself, or absolutely.
				}
		saying: I have overcome and have trodden the wine-press alone,%
			\footnote{%
				% \label{}%
				This language comes from \hyperref[ISAIAH-63-3]{Isaiah 63:3}. See also \hyperref[DC88-106]{DC 88:106}.
				}
		even the wine-press of the fierceness of the wrath of Almighty God.

		% \begin{framed}
		% \scriptsize
			% \begin{compactdesc}
				% \item[JSP] 
				% % \item[BC] 
				% % \item[]
			% \end{compactdesc}
		% \end{framed}

	% -----
	\subsubsection{Section 76:108} % *DC76-108 
	% -----
		\label{DC76-108}

		Then shall he be crowned with the crown of his glory, to sit on the throne of his power to reign forever and ever.

		% \begin{framed}
		% \scriptsize
			% \begin{compactdesc}
				% \item[JSP] 
				% % \item[BC] 
				% % \item[]
			% \end{compactdesc}
		% \end{framed}

	% -----
	\subsubsection{Section 76:109} % *DC76-109 
	% -----
		\label{DC76-109}

		But behold, and lo, we saw the glory and the inhabitants of the telestial world, that they were as innumerable as the stars in the firmament of heaven, or as the sand upon the seashore;

		% \begin{framed}
		% \scriptsize
			% \begin{compactdesc}
				% \item[JSP] 
				% % \item[BC] 
				% % \item[]
			% \end{compactdesc}
		% \end{framed}

	% -----
	\subsubsection{Section 76:110} % *DC76-110 
	% -----
		\label{DC76-110}

		And heard the voice of the Lord saying: These all shall bow the knee, and every tongue shall confess to him who sits upon the throne forever and ever;

		% \begin{framed}
		% \scriptsize
			% \begin{compactdesc}
				% \item[JSP] 
				% % \item[BC] 
				% % \item[]
			% \end{compactdesc}
		% \end{framed}

	% -----
	\subsubsection{Section 76:111} % *DC76-111 
	% -----
		\label{DC76-111}

		For they shall be judged according to their works, and every man shall receive according to his own works, his own dominion,%
			\footnote{%
				% \label{}%
				Note that in the JSP version, the EMS, and the 1835 and 1844 DC, you get an ``and'' before the ``his own dominion'' bit. HDDC 970 notes that this ``and'' was present in many early versions beyond just those I've listed below. Taking away the ``and'' makes a pretty serious difference in the verse. As it currently appears, it's kind of an awkward sentence, but the most straightforward way of taking it is to see ``his own dominion'' as the object of ``receive.'' We could punctuate it in the following way to bring out this interpretation more clearly:
				
					\begin{tabular}{p{2cm}p{15cm}}
						\textbf{No ``and''} & For they shall be judged according to their works, \textbf{and every man shall receive, according to his own works, his own dominion in the mansions which are prepared}.
					\end{tabular}
					
				This makes sense---but this is not what the punctuation is in the current version, so it isn't even clear that this is the right way to understand it. I like adding back in the ``and,'' since it produces a more interesting idea. Here is a way to do that, with punctuation and other aids to suggest an interpretation:
				
					\begin{tabular}{p{2cm}p{15cm}}
						\textbf{With ``and''} & For they shall be judged according to their works. And every man shall receive (i) according to his own works, and (ii) [according to] his own dominion, in the mansions which are prepared.
					\end{tabular}
					
				On this interpretation, ``receive'' is intransitive and so doesn't take an object. The interesting thing that this interpretation adds---beyond being true to the original text---is that there is a new condition added that determines how we receive, and that condition is that we receive according to our ``own dominion.''
				}
		in the mansions which are prepared;

		\begin{framed}
		\scriptsize
			\begin{compactdesc}
				\item[JSP] for they shall be judged according to there works and every man shall receive according to his own works and his own dominion in the mansions which are prepared,
				\item[EMS] for they shall be judged according to their works; and every man shall receive according to his own works, and his own dominion, in the mansions which are prepared;
				\item[BC] ---
				\item[DC 1835] for they shall be judged according to their works; and every man shall receive according to his own works, and his own dominion, in the mansions which are prepared,
				\item[DC 1844] for they shall be judged according to their works; and every man shall receive according to his own works, and his own dominion, in the mansions which are prepared,
			\end{compactdesc}
		\end{framed}

	% -----
	\subsubsection{Section 76:112} % *DC76-112 
	% -----
		\label{DC76-112}

		And they shall be servants of the Most High; but where God and Christ dwell they cannot come, worlds without end.

		% \begin{framed}
		% \scriptsize
			% \begin{compactdesc}
				% \item[JSP] 
				% % \item[BC] 
				% % \item[]
			% \end{compactdesc}
		% \end{framed}

	% -----
	\subsubsection{Section 76:113} % *DC76-113 
	% -----
		\label{DC76-113}

		This is the end of the vision which we saw, which we were commanded to write while we were yet in the Spirit.

		% \begin{framed}
		% \scriptsize
			% \begin{compactdesc}
				% \item[JSP] 
				% % \item[BC] 
				% % \item[]
			% \end{compactdesc}
		% \end{framed}

	% -----
	\subsubsection{Section 76:114} % *DC76-114 
	% -----
		\label{DC76-114}

		But great and marvelous are the works of the Lord, and the mysteries of his kingdom which he showed unto us, which surpass all understanding in glory, and in might, and in dominion;

		% \begin{framed}
		% \scriptsize
			% \begin{compactdesc}
				% \item[JSP] 
				% % \item[BC] 
				% % \item[]
			% \end{compactdesc}
		% \end{framed}

	% -----
	\subsubsection{Section 76:115} % *DC76-115 
	% -----
		\label{DC76-115}

		Which he commanded us we should not write while we were yet in the Spirit, and are not lawful for man to utter;

		% \begin{framed}
		% \scriptsize
			% \begin{compactdesc}
				% \item[JSP] 
				% % \item[BC] 
				% % \item[]
			% \end{compactdesc}
		% \end{framed}

	% -----
	\subsubsection{Section 76:116} % *DC76-116 
	% -----
		\label{DC76-116}

		Neither is man capable to make them known, for they are only to be seen and understood by the power of the Holy Spirit, which God bestows on those who love him, and purify themselves before him;

		% \begin{framed}
		% \scriptsize
			% \begin{compactdesc}
				% \item[JSP] 
				% % \item[BC] 
				% % \item[]
			% \end{compactdesc}
		% \end{framed}

	% -----
	\subsubsection{Section 76:117} % *DC76-117 
	% -----
		\label{DC76-117}

		To whom he grants this privilege of seeing and knowing for themselves;

		% \begin{framed}
		% \scriptsize
			% \begin{compactdesc}
				% \item[JSP] 
				% % \item[BC] 
				% % \item[]
			% \end{compactdesc}
		% \end{framed}

	% -----
	\subsubsection{Section 76:118} % *DC76-118 
	% -----
		\label{DC76-118}

		That through the power and manifestation of the Spirit, while in the flesh, they may be able to bear his presence in the world of glory.

		% \begin{framed}
		% \scriptsize
			% \begin{compactdesc}
				% \item[JSP] 
				% % \item[BC] 
				% % \item[]
			% \end{compactdesc}
		% \end{framed}

	% -----
	\subsubsection{Section 76:119} % *DC76-119 
	% -----
		\label{DC76-119}

		And to God and the Lamb be glory, and honor, and dominion forever and ever. Amen.

		% \begin{framed}
		% \scriptsize
			% \begin{compactdesc}
				% \item[JSP] 
				% % \item[BC] 
				% % \item[]
			% \end{compactdesc}
		% \end{framed}